\section{Governance: decentralization, operator protection, and the content-reporting quorum}
\label{sec:23-governance}\label{sec:governance}

Every network that executes other people's software eventually has to answer a question that
storage networks can defer and payment networks can ignore: who may stop a workload from running,
and by what procedure. Flux answers it twice. There is an answer that is deployed — a small set of
JSON documents in a GitHub repository that every node fetches on a timer and enforces without
argument — and there is an answer that is designed but not built: a collateral-weighted vote,
conducted entirely in chain transactions, in which a node operator reports an application and an
absolute fraction of total voting weight decides whether it is banned. This section states the
deployed answer as a mechanism rather than as an implementation detail, specifies the designed one
at the depth a reimplementer would need, proves what can be proved about it, and then reports the
measurement that decides its worth.

The single claim the section carries is this: \textbf{collateral weighting makes the quorum
Sybil-resistant by construction and, at the distribution measured on 2026-09-03, not
censorship-resistant.} Four owner identities hold enough collateral-weighted voting power to evict any
application by themselves --- three, once addresses linked by a shared node key or a shared collateral
funding transaction are merged. That is a serious finding about
the proposal. It is also, and simultaneously, an improvement on what is deployed, where the
corresponding number is one, the decision carries no signature, and no node can verify that it was
ever made. Both statements are true and the section makes both.

A second property of the specified design is worth stating in the same breath, because it is what
makes the first one auditable: \textbf{nothing is tallied off-chain.} The report is a transaction,
each vote is a transaction, the count is a function \fluxd{} evaluates from those transactions, and
the ban is a memo in a coinbase. Every step is a chain fact that any node can recompute from block
data alone, with no tallier, no publisher key and no API in the trust path.

\begin{remark}[Local notation]
This section uses, beyond \texttt{notation.tex}: $\vtotal^{\!\star}$ for the trailing-maximum total
weight, $\kappa(\vthresh)$ for the size of the smallest colluding set that reaches threshold
$\vthresh$, $Y_r$ for a report's tally, $\mathrm{own}_j$ for a node's \emph{owner} (collateral) key
and $\mathrm{pk}_j$ for its VPS message-signing key --- these are different keys and only the first
is an identity here (\Cref{def:23:electorate}) --- $X_i$ for an identity's proposal capacity, $s_i$
for a node's share of the \PNR pool, $m$ and $T$ for the identity-chamber count and tenure, and
$\Delta_{\mathrm{vote}}, \Delta_{\mathrm{exec}}, \Delta_{W}$ for windows measured in blocks. All are
proposed additions to the notation file; none overloads an existing symbol. Note in particular that
$\kappa$ denotes a \emph{collusion-set size} throughout and never a key. The report bond $\bond$
declared in \texttt{notation.tex} is \emph{not} used: the specified design has no bond, and prices
reporting by collateral instead (\Cref{def:23:capacity}).
\end{remark}

% =====================================================================
\subsection{Governance as it operates today}
\label{sec:23:today}

Flux already has content governance. It is not a quorum, it is not on-chain, and it is not
described in the public roadmap. It is a publisher.

\begin{definition}[Deployed network-policy channel, \shipped]
\label{def:23:policy-channel}
The channel is the tuple $\bigl(\mathcal{D}, \mathcal{O}_{\mathrm{pub}}, \mathrm{fetch},
\mathrm{enforce}\bigr)$ where:
\begin{itemize}
  \item $\mathcal{D}$ is the document set $\{$\texttt{blockedrepositories.\allowbreak{}json},
    \texttt{vettedrepositories.\allowbreak{}json}, \texttt{enterprisenodes.\allowbreak{}json},
    \texttt{tamperingblockednodes.\allowbreak{}json}, \texttt{iplocation.\allowbreak{}bin.\allowbreak{}gz}$\}$, served as unsigned
    documents over plain HTTPS from a GitHub repository
    \srccode{fluxos-network-policy/\allowbreak{}README.md}{1--9};
  \item $\mathcal{O}_{\mathrm{pub}}$ is the publisher set: the owners of the \texttt{RunOnFlux}
    GitHub organisation, of whom the repository states ``everyone who has it is an admin'', with
    branch protection forbidding direct pushes to \texttt{main} but \emph{deliberately} no
    second-reviewer requirement — the stated reasoning being that a rule requiring a second person
    ``gets bypassed at 3am'' \srccode{fluxos-network-policy/\allowbreak{}.github/\allowbreak{}CODEOWNERS}{1--15}
    \srccode{fluxos-network-policy/\allowbreak{}README.md}{297--311};
  \item $\mathrm{fetch}$ is a per-node scheduled retrieval — \SI{6}{\hour} for
    \texttt{blockedrepositories.\allowbreak{}json} and \texttt{enterprisenodes.\allowbreak{}json}, \SI{12}{\hour} for
    \texttt{tamperingblockednodes.\allowbreak{}json} \srccode{fluxos-network-policy/\allowbreak{}README.md}{30--32}, and ---
    on the \specv{9} branch only --- a daily conditional fetch for the geolocation table
    \srcbranch{MorningLightMountain713/\allowbreak{}flux}{feat/fluxos-v9}{ZelBack/config/default.js:779--784} — followed by shape validation
    (\texttt{isStringArray}, \texttt{isNodeOwnerMap}; row/country/org/region floors for the
    geolocation artifact) \srccode{fluxos-network-policy/\allowbreak{}scripts/\allowbreak{}validate.js}{22--60};
  \item $\mathrm{enforce}$ is the per-node compliance sweep: \FluxOS{} fetches
    \texttt{helpers/\allowbreak{}blockedrepositories.\allowbreak{}json} from the raw-content URL of \texttt{RunOnFlux/flux},
    caches it \srccode{flux/\allowbreak{}ZelBack/\allowbreak{}src/\allowbreak{}services/\allowbreak{}appSecurity/\allowbreak{}imageManager.js}{179--195}, and every
    \texttt{image\allowbreak{}Compliance\allowbreak{}Interval\allowbreak{}Ms} $=$ \SI{1}{\hour}
    \srccode{flux/\allowbreak{}ZelBack/\allowbreak{}config/\allowbreak{}default.js}{319}
    \srccode{flux/\allowbreak{}ZelBack/\allowbreak{}src/\allowbreak{}services/\allowbreak{}serviceManager.js}{529--531} removes every locally installed
    application whose image, owner address or specification hash matches, spacing removals
    \SI{3}{\minute} apart
    \srccode{flux/\allowbreak{}ZelBack/\allowbreak{}src/\allowbreak{}services/\allowbreak{}appSecurity/\allowbreak{}imageManager.js}{644--676}. A parallel
    \texttt{vettedrepositories.\allowbreak{}json} exempts listed applications from blocked ports and blocked
    repositories \srccode{flux/\allowbreak{}ZelBack/\allowbreak{}src/\allowbreak{}services/\allowbreak{}appSecurity/\allowbreak{}imageManager.js}{202--236,395--410}.
\end{itemize}
The node-local state machine is fail-frozen. A fetch that fails or fails shape validation leaves
the previous copy in force indefinitely and across restarts; an empty document (\texttt{[]}) is
honoured as ``nothing is blocked'' and is explicitly distinguished from an unreachable one; a node
that has never fetched successfully declines to answer rather than guessing
\srccode{fluxos-network-policy/\allowbreak{}README.md}{33--44}.
\end{definition}

\begin{algorithm}[H]
\caption{Deployed policy fetch-and-enforce loop, per node (\shipped)}
\KwIn{document set $\mathcal{D}$; publisher URL; per-document fetch interval; local application set $A_{\mathrm{loc}}$}
\KwOut{$A_{\mathrm{loc}}$ with policy-matching applications removed}
\tcc{Invariant: a node never enforces a document it has not itself shape-validated, and never
downgrades to ``nothing blocked'' because of a transport failure.}
\tcc{Failure mode: fetch failure freezes the last valid copy indefinitely (stale enforcement, and
stale geolocation for placement); never-fetched nodes decline to answer. There is no signature to
check, so a MITM or a compromised publisher account is indistinguishable from a legitimate update.}
\ForEach{$d \in \mathcal{D}$ due for refresh}{
  $x \leftarrow$ \texttt{HTTPS GET}$(d)$\;
  \lIf{transport error}{keep cached $d$; \textbf{continue}}
  \lIf{$\neg\,\mathrm{shapeValid}(x)$}{reject $x$; keep cached $d$; \textbf{continue}}
  cache $x$ as the current copy of $d$\;
}
\ForEach{application $\app \in A_{\mathrm{loc}}$ (hourly sweep)}{
  \If{image, owner address or specification hash of $\app$ matches \texttt{blockedrepositories}
      \textbf{and} $\app$ is not in \texttt{vettedrepositories}}{
    remove $\app$ locally; wait \SI{3}{\minute} before considering the next removal\;
  }
}
\end{algorithm}

A change to any of these documents takes effect on every node at its next scheduled fetch — no
staged rollout, no canary, no version negotiation \srccode{fluxos-network-policy/\allowbreak{}README.md}{3--5}.
What such a change can do is broad and precise: block any container image or namespace, any
application owner address, or any application hash network-wide; grant or revoke which owner
addresses may install on which enterprise node; mark a node's collateral transaction hash as DOSed
for tampering; exempt an owner, hash or repository from user-level blocks; and, through the monthly
baseline pull request, change which organisations and regions are treated as residential rather
than hosting for placement and enforcement purposes
\srccode{fluxos-network-policy/\allowbreak{}README.md}{99--196,320--326}. The documents are unsigned; the
repository states the position plainly — ``their authenticity rests on GitHub and on who can merge
to this repo'' — and records signed documents with a rollback-resistant sequence number as the
intended next step rather than a finished one \srccode{fluxos-network-policy/\allowbreak{}README.md}{320--326}.

A second lever sits in the reachability plane. The Foundation-operated \FDM{} fleet carries a
configuration-level wildcard blacklist and whitelist by application name, plus an
\texttt{ownersApps} restriction that can confine an entire \FDM{} instance to one owner's
applications \srccode{flux-domain-manager/\allowbreak{}config/\allowbreak{}appsConfig.js}{2--6}
\srccode{flux-domain-manager/\allowbreak{}src/\allowbreak{}services/\allowbreak{}application/\allowbreak{}subset.js}{25--47}, and a hardcoded
\texttt{forbidden-\allowbreak{}backend} hostname list shipped to every \FDM{} instance in the HAProxy template
\srccode{flux-domain-manager/\allowbreak{}src/\allowbreak{}services/\allowbreak{}haproxyTemplate.js}{499--512}. This lever is exercised:
the commit history contains ``forbid domains'' and three ``ban phishing website'' changes. A third
sits at consensus level: a hardcoded 2-of-4 key set can produce blocks outside normal consensus,
with no time gating and no stall detection, despite a source comment saying gating was planned
\srccode{fluxd/\allowbreak{}src/\allowbreak{}emergencyblock.cpp}{24--30,183--190}
\srccode{fluxd/\allowbreak{}src/\allowbreak{}emergencyblock.h}{40--44}; the deployed threshold is
\texttt{n\allowbreak{}Emergency\allowbreak{}Min\allowbreak{}Signatures} $=2$ against a four-key set
\srccode{fluxd/\allowbreak{}src/\allowbreak{}chainparams.cpp}{253,522}, and the team states 3-of-4 with a newly generated key
set as the intended configuration, aligning consensus emergency-key governance with the bridge's
3-of-4 quorum (\S22), with no activation scheduled
\srcintent{evidence/design/08-answers-round-2.md, round 5 answers} (\planned). That key set is not a
content lever, but it is a governance lever, and a section on who
can do what to this network must list it. Under the mechanism specified below it acquires a second
significance, because that mechanism's enforcement rides in the coinbase
(\Cref{sec:23:enforcement}).

\begin{assumption}[Trust boundaries invoked here]
\label{ass:23:today}
As fixed in the system-model section: node operators are mutually distrusting and independently
operated; the \FluxOS{} message layer is a gossip mesh with per-peer bad-message accounting
(five malformed or badly signed messages in ten minutes closes the connection)
\srccode{flux/\allowbreak{}ZelBack/\allowbreak{}src/\allowbreak{}services/\allowbreak{}fluxCommunication.js}{597--618}; \fluxd's deterministic node set
is consensus state; and the policy channel, the \FDM{} fleet and the application-discovery API are
Flux-Foundation-operated. Nothing in \Cref{def:23:policy-channel} is consensus state, and no
transition of it is verifiable by a node against the chain.
\end{assumption}

\begin{property}[The collusion measure of the deployed mechanism, \shipped]
\label{prop:23:kappa-today}
For the channel of \Cref{def:23:policy-channel}: the number of parties whose agreement is required
to remove an arbitrary application from every node is $\kappa_{\mathrm{deployed}} = 1$; the
authentication on the decision is $\varnothing$ (TLS to a code-hosting provider, no document
signature); the record available to a node for verifying that the decision was authorised is
$\varnothing$; and the latency from decision to fleet-wide effect is bounded by the fetch interval
plus the sweep interval, i.e.\ within \SI{7}{\hour} for the repository blocklist and \SI{13}{\hour}
for the tampering list.
\end{property}

\begin{proof}[Proof sketch]
A single organisation owner opens and merges a pull request adding an image, owner address or
application hash to \texttt{blockedrepositories.\allowbreak{}json}; CI checks shape only and no second reviewer
is required \srccode{fluxos-network-policy/\allowbreak{}README.md}{297--308}. Every node fetches on its own
schedule and enforces the new document without any authentication step, because none exists
\srccode{fluxos-network-policy/\allowbreak{}README.md}{320--326}. The bound on latency follows from the
\SI{6}{\hour}/\SI{12}{\hour} fetch intervals and the \SI{1}{\hour} sweep. That the record is not
node-verifiable follows from the absence of a signature: a node observing the document cannot
distinguish a legitimate merge from a substituted document.
\end{proof}

This is the baseline against which the quorum must be judged, and it is the strongest argument for
building the quorum at all. The proposed mechanism, with every defect catalogued below, would
require four parties rather than one --- three under the strongest linkage computable from public
data --- would leave a permanent chain record naming exactly which keys produced the eviction, and
would freeze its electorate at a chain height. It would be a large
decentralization improvement over what is deployed. It would still not be censorship-resistant.
The section says both, in that order, because presenting the quorum as a solution to a problem the
network does not have would be the more comfortable and the less honest framing.

% =====================================================================
\subsection{The roadmap principle, and the tension the quorum creates}
\label{sec:23:c5}

The public roadmap does not contain a content-eviction vote. What it contains is a principle
pointing the other way:

\begin{quote}
``The principle: defence as protection from harm, not veto over content. \textbf{Behaviour, not
content.}'' \srcroadmap{flux-roadmap-public.md:94}
\end{quote}

Under the heading of operator protection the roadmap lists an application oracle, resource
guardrails, an optional egress profile, and an evidence pack — operator-side defences that let a
node refuse or contain a workload. It does not list a network vote that removes an application from
every node. The mechanism specified below is therefore not an elaboration of a published
commitment; it is a design-intake item \srcintent{evidence/design/04-moderation.md} that must be
reconciled with a published principle, and this paper does the reconciling in the open rather than
letting a reader find the gap.

The reconciliation has one hinge, and it is an open parameter: the category taxonomy. If the
reportable categories are behavioural — harm to third parties, harm to operators, harm to the
network's own infrastructure — then the quorum is an enforcement path for the roadmap's own
principle and the two documents agree. If the categories are about what an application says or
serves, the quorum is a content veto and the two documents contradict each other. Nothing in the
mechanism forces either reading; the taxonomy does, and while its enumeration is still undecided,
its scope has since been ruled on (\Cref{rem:content-veto}).

\begin{remark}[The scope is a network-level content veto --- settled, and it supersedes the roadmap]
\label{rem:content-veto}
This paper asked directly whether the report taxonomy should be restricted to behavioural categories
--- harm to third parties, to operators, to the network --- so as to reconcile the quorum with the
roadmap's published ``behaviour, not content'' principle, noting that any content-referring category
turns the quorum into a network-level content veto. The answer is that it is a network-level veto
\srcintent{08-answers-round-2.md, round 9} (\planned).

The paper records this as a deliberate decision and states its consequences without softening them,
because they are the consequences that matter to anyone evaluating Flux as infrastructure.

\begin{enumerate}
  \item \textbf{It supersedes a published position.} The roadmap's ``behaviour, not content'' framing
  \srcroadmap{flux-roadmap-public.md} described a narrower mechanism than the one now intended. Where
  this paper previously reasoned from that principle --- including in
  \Cref{sec:23:priorart}'s comparison with IPFS, Filecoin and Arweave --- the comparison still
  holds on the facts but no longer on the stated intent: Flux is choosing a stronger power than those
  systems exercise, deliberately.
  \item \textbf{It makes concentration the load-bearing property.} A behavioural taxonomy limits the
  damage a captured quorum can do, because the categories themselves resist misuse. A content veto
  has no such internal limit: what the quorum may remove is bounded only by what the quorum will
  vote for. The measurement in \Cref{sec:23:measurement} --- four owner identities reaching the
  \SI{25}{\percent} threshold, three of them under the strongest linkage computable from public data
  --- therefore stops being a background statistic and becomes the mechanism's principal risk.
  \item \textbf{It does not change the deterrence analysis.} Removal plus forfeiture of a prepaid
  balance still has essentially no deterrent value against a determined deployer who redeploys under
  a new name; the author confirms there is no reinstatement and that redeployment as a fresh
  application is expected \srcintent{08-answers-round-2.md, round 5}.
\end{enumerate}

What remains genuinely open is not the scope but the enumeration: which categories exist, so that a
report is a well-formed object with a fixed domain rather than free text. That is a specification
task, not a policy question, and \Cref{sec:23:mechanism} states the requirement it must satisfy.
\end{remark}

There is a second, quieter tension. The mechanism's enforcement is removal, and its financial
consequence for the deployer is forfeiture of an unused prepaid balance whose floor is
\num{0.01}\flux\ (\Cref{sec:23:enforcement}). The purpose is to stop harm, not to fine, and under
either scope that is the right design. It also means the mechanism has essentially no
deterrent value against a determined deployer, who redeploys under a new name for the price of one
transaction. The paper states this as a property of the design rather than as a criticism of it,
because it follows directly from the principle the design is trying to honour.

% =====================================================================
\subsection{The content-reporting quorum, specified}
\label{sec:23:mechanism}

Everything in this subsection is \planned. No part of it exists in any repository surveyed for this
paper; the settled inputs are the design intake \srcintent{evidence/design/04-moderation.md}, the
on-chain construction supplied by the team
\srcintent{evidence/design/08-answers-round-2.md} and the measured node set
\srcmeasured{api.runonflux.io/\allowbreak{}daemon/\allowbreak{}listzelnodes}{2026-09-03}. Settled and not re-litigated here:
collateral weighting at one unit per \num{1000}\flux; votes signed with node owner keys;
$\vthresh = 0.25$ as an \emph{absolute} fraction of total weight; abstention counts as ``no'';
enforcement is cancellation and removal with forfeiture of the tenant's prepaid balance; node
collateral is never slashed.

Settled additionally, and this is the part that decides the section's verifiability requirement:
the report is an on-chain transaction shaped after the deployed \FluxNode{} start and confirm
types, carrying the application hash to be banned and indexed so that any node can find it; a
\SI{24}{\hour} window (\num{2880} blocks at \SI{30}{\second} spacing) opens on its confirmation;
each vote is a second special transaction; \fluxd{} counts the quorum internally from those
transactions; on the threshold being reached the block template changes and the \PoN{} producer
carries a ban memo in the coinbase, at most one per block; \FluxOS{} observes the ban and purges
the application; vote transactions are fee-less; an unmet report expires and the application may be
re-reported immediately; and a key may hold at most as many concurrently open proposals as it
operates nodes \srcintent{evidence/design/08-answers-round-2.md}.

One consequence of that shape should be stated before the construction rather than after it. This
design puts moderation into \emph{consensus}. It needs two new transaction types, a new coinbase
rule and new chain-derived state in \fluxd, which is a network upgrade; the message-layer
alternative it replaces needed only a \FluxOS{} release. The consensus route buys exact
determinism and independent verifiability (\Cref{thm:23:verifiable}) and pays for them in
deployment cost, and this network's own history of what a consensus change costs --- the \PoN{}
activation required a temporary deep-reorg allowance and a cluster of emergency checkpoints
\srccode{fluxd/\allowbreak{}src/\allowbreak{}main.cpp}{8983--8988} \srccode{fluxd/\allowbreak{}src/\allowbreak{}chainparams.cpp}{277--283} --- is the
right reference class for estimating it.

\subsubsection{Electorate and weight}

\begin{definition}[Electorate, identity, weight, reference total, \planned]
\label{def:23:electorate}
Let $N(\hgt)$ be the deterministic confirmed \FluxNode{} set at height $\hgt$, which is consensus
state in \fluxd. Each node $j \in N(\hgt)$ carries a tier $t_j \in \tiers$, collateral
$\coll{t_j} \in \{\num{1000},\ \num{12500},\ \num{40000}\}\flux$, and \emph{two distinct keys}: the
owner key $\mathrm{own}_j$, which is the collateral public key and from which the record's
\texttt{payment\_address} is derived
\srccode{fluxd/\allowbreak{}src/\allowbreak{}rpc/\allowbreak{}fluxnode.cpp}{1512--1520}, and the node key $\mathrm{pk}_j$, which
\fluxd{} documents as ``\texttt{Pubkey used for VPS signature verification}''
\srccode{fluxd/\allowbreak{}src/\allowbreak{}fluxnode/\allowbreak{}fluxnode.h}{154} and which signs \PoN{} blocks and \FluxOS{} P2P
envelopes. Because the settled rule is that votes are signed with node \emph{owner} keys, an
\emph{identity} here is an owner key and never a node key; the identity set is
$I(\hgt) = \{\mathrm{own}_j : j \in N(\hgt)\}$, publicly proxied by \texttt{payment\_address}, which
is the observable the measurement of \Cref{sec:23:measurement} uses. Define
$\vweight{\cdot} : I(\hgt) \to \mathbb{R}_{\ge 0}$ and $\vtotal : \mathbb{N} \to \mathbb{R}_{>0}$ by
\[
\vweight{i}(\hgt) \;\defeq\; \sum_{j \in N(\hgt):\ \mathrm{own}_j = i}
  \frac{\coll{t_j}}{\num{1000}\flux},
\qquad
\vtotal(\hgt) \;\defeq\; \sum_{i \in I(\hgt)} \vweight{i}(\hgt)
  \;=\; \sum_{t \in \tiers} \ncount{t}(\hgt)\,\frac{\coll{t}}{\num{1000}\flux}.
\]
Per node the weights are $\tC \mapsto 1$, $\tN \mapsto 12.5$, $\tS \mapsto 40$. The
\emph{reference total} is a trailing maximum,
\[
\vtotal^{\!\star}(\hgt_0) \;\defeq\; \max_{\hgt \in [\hgt_0 - \Delta_W,\ \hgt_0]} \vtotal(\hgt),
\qquad \Delta_W \in \mathbb{N}\ \text{open}.
\]
All weights in a report are evaluated at the height $\hgt_0$ at which the report transaction
confirms and never re-evaluated: a node confirmed after $\hgt_0$ would carry weight $0$ in that
report.
\end{definition}

Two design points deserve their reasons. The \emph{electorate freeze} exists because collateral
maturity is 100 blocks (\SI{50}{\minute} at \SI{30}{\second} spacing)
\srccode{fluxd/\allowbreak{}src/\allowbreak{}fluxnode/\allowbreak{}fluxnode.h}{20}; without the freeze a report
could be answered by nodes created in response to it. The \emph{trailing maximum} exists because
the reporter chooses $\hgt_0$: node confirmations expire after 640 blocks (\SI{5.3}{\hour}), so an
incident that leaves a large share of weight temporarily unconfirmed would lower $\vtotal(\hgt_0)$
and raise every remaining share. Taking a maximum rather than a mean is the safe direction — it can
only make eviction harder.

\settlednote{$\Delta_W = \num{20160}$ blocks, as a trailing maximum (\Cref{tab:23-recommended}). $\Delta_W$, the trailing window for $\vtotal^{\!\star}$. Domain: $[0,\ \num{20160}]$
blocks. Trade-off: $\Delta_W = 0$ reduces to the plain total and admits the electorate-timing attack
of \Cref{sec:23:attacks}; large values leave $\vtotal^{\!\star}$ stale after a genuine fleet
contraction, which raises the bar to evict and is therefore safe but costs liveness.}

At the measurement height \num{2917115}: $\ncount{\tC} = \num{3247}$, $\ncount{\tN} = \num{1615}$,
$\ncount{\tS} = \num{1627}$, so $\vtotal = \num{3247} + \num{20187.5} + \num{65080} = \num{88514.5}$
units, i.e.\ \num{88514500}\flux\ of locked collateral, with tier shares
\SI{3.67}{\percent} / \SI{22.81}{\percent} / \SI{73.52}{\percent}
\srcmeasured{api.runonflux.io/\allowbreak{}daemon/\allowbreak{}listzelnodes}{2026-09-03}.

\subsubsection{Three transaction types}

\begin{definition}[Report transaction, \planned]
\label{def:23:report}
A report is a \fluxd{} transaction of a new \FluxNode-family type $\mathsf{MOD\_REPORT}$, carrying
\[
r \;=\; \bigl(\mathsf{app\_hash},\ \mathsf{category},\ \mathsf{evidence},\ \mathrm{own}_r\bigr),
\qquad
\sigma_r \;=\; \mathrm{Sign}_{\mathrm{own}_r}\bigl(H(r)\bigr),
\]
where $\mathsf{app\_hash} \in \{0,1\}^{256}$ is the application hash to be banned, $\mathsf{category}$
is drawn from the taxonomy left open above, $\mathsf{evidence}$ is a content hash of supporting
material, and $\mathrm{own}_r \in I(\hgt_0)$ is the reporter's owner key. The report identifier is
$H(r)$; the transaction is indexed by $\mathsf{app\_hash}$ and by $H(r)$, so that any node can
enumerate the open reports on an application from its own chain state. It carries no fee and no
bond, and it is valid only if no report on the same $\mathsf{app\_hash}$ is open at its confirmation
height (\Cref{def:23:decision}) and $\mathsf{app\_hash}$ is not already in the ban set
(\Cref{def:23:ban}) \srcintent{evidence/design/08-answers-round-2.md, round 18}. Its confirmation
height defines $\hgt_0$ and opens the window
$[\hgt_0,\ \hgt_0 + \Delta_{\mathrm{vote}}]$ with $\Delta_{\mathrm{vote}} = \num{2880}$ blocks,
\SI{24}{\hour} at \SI{30}{\second} spacing.
\end{definition}

The transaction shape is not new: \fluxd{} already carries two special \FluxNode{} types --- the
start transaction, signed by the collateral key
\srccode{fluxd/\allowbreak{}src/\allowbreak{}fluxnode/\allowbreak{}activefluxnode.cpp}{253--288}, and the confirm transaction, signed by
the node key and then benchmark-signed \srccode{fluxd/\allowbreak{}src/\allowbreak{}fluxnode/\allowbreak{}activefluxnode.cpp}{290--332,59--68} ---
carried at transaction versions 5 and 6
\srccode{fluxd/\allowbreak{}src/\allowbreak{}fluxnode/\allowbreak{}activefluxnode.cpp}{409--427}. The report and vote types would be a
third and fourth member of that family.

Two boundaries of what a chain can check follow immediately, and both are properties of the design
rather than defects in it. First, \fluxd{} holds no application registry: the \FluxOS{} global
specification set is not consensus state. Consensus can therefore validate a report's \emph{form}
(type, signature, capacity) and never its \emph{referent}; a report naming a hash that corresponds
to no deployed application is perfectly valid on-chain and simply unenforceable off it. Second,
$\reporters$ --- reporter eligibility, an open parameter in every earlier draft of this mechanism ---
is settled by the anti-spam rule below rather than chosen: a key that operates no nodes has
capacity $0$, so $\reporters = I(\hgt_0)$ falls out of \Cref{def:23:capacity}.

\settlednote{$(\mathsf{app\_name}, \mathsf{owner})$ (\Cref{tab:23-recommended}). What $\mathsf{app\_hash}$ hashes, which is also the enforcement scope. Domain:
$\{$the permanent-message hash of one specification version; a hash of the application identity
$(\mathsf{app\_name}, \mathsf{owner})$; $+$ all applications of the same owner; $+$ image$\}$.
Trade-off: if the ban key is a specification-version hash, an application updated during the
\SI{24}{\hour} window acquires a new hash and the ban, when it lands, names a version no longer
deployed --- the escape that keying on name-and-owner was chosen to close, and the invariant
``an update during the window is neither an escape nor a second claim'' holds only under the
identity reading. Owner scope is evaded by registering a new address at zero cost while widening
the blast radius of a mistaken ban; image scope is what the deployed blocklist does and is defeated
by re-tagging. The team's specification says ``the application hash'' and does not disambiguate.}

\begin{definition}[Vote transaction, \planned]
\label{def:23:vote}
A vote is a \fluxd{} transaction of type $\mathsf{MOD\_VOTE}$ carrying
\[
v \;=\; \bigl(H(r),\ \hgt_0,\ \mathsf{app\_hash},\ i,\ \mathsf{decision}\bigr),
\qquad
\sigma_v \;=\; \mathrm{Sign}_{i}\bigl(H(v)\bigr),
\]
with $i \in I(\hgt_0)$ and $\mathsf{decision} \in \{\mathsf{evict}, \mathsf{keep}\}$. It carries no
fee. It is valid iff its signer is in the frozen electorate $I(\hgt_0)$, its confirmation height
$\hgt_v$ satisfies $\hgt_0 \le \hgt_v \le \hgt_0 + \Delta_{\mathrm{vote}}$, and no earlier vote
from $i$ on $H(r)$ is already in the chain. Votes are therefore \emph{irrevocable within a report}:
the first confirmed vote from an identity is the only one that counts, and a later one is invalid
rather than overriding.
\end{definition}

What is signed binds the vote to one report and one electorate, so it cannot be replayed against a
later report on the same application (different $\hgt_0$) or against a re-registration. Two things
change relative to the message-layer construction this replaces, and both are improvements. The
window is measured on \emph{confirmation} height, which is a canonical clock rather than a
self-declared envelope timestamp, so two honest nodes cannot disagree about whether a vote arrived
in time --- at the cost that a vote broadcast near the boundary and confirmed after it is simply
void, which the \SI{24}{\hour} window makes an edge case rather than a hazard. And irrevocability
is now a consequence of chain order rather than a rule imposed to make an off-chain certificate
checkable. Under the settled rule $\mathsf{keep}$ votes are never counted.

\settlednote{Evict-only, with a per-identity mempool quota as the relay policy (\Cref{tab:23-recommended}). Whether the vote type carries a $\mathsf{decision}$ field at all, and the relay policy for
fee-less transactions. Domains: $\mathsf{decision} \in \{$evict-only; evict/keep with keep recorded
and uncounted$\}$; relay policy $\in \{$admit up to the capacity bound; per-identity mempool
quota; \ldots$\}$. Trade-off: under abstention-as-no a $\mathsf{keep}$ vote changes no outcome and
costs block space, so recording one buys only a public record and an input to a possible
early-termination rule; and because these transactions pay no fee there is no price signal to order
a mempool by, so admission must be bounded by the capacity rule and by one-vote-per-identity
instead. A node that drops fee-less transactions under pressure delays a vote but cannot invalidate
it, since 2{,}880 slots remain in the window.}

\begin{definition}[Tally, decision rule, expiry, \planned]
\label{def:23:decision}
Let $V_r(\hgt) \subseteq I(\hgt_0)$ be the set of identities whose valid $\mathsf{evict}$ vote on
$r$ is \emph{confirmed at or before} height $\hgt$, and let
$Y_r : \mathbb{N} \to \mathbb{R}_{\ge 0}$, $Y_r(\hgt) \defeq \sum_{i \in V_r(\hgt)}
\vweight{i}(\hgt_0)$. Every validator evaluates $Y_r$ from block data; there is no tallier. The
report is \emph{carried} iff
\[
\mathsf{carried}(r) \iff \exists\, \hgt_c \in [\hgt_0,\ \hgt_0 + \Delta_{\mathrm{vote}}]
:\quad Y_r(\hgt_c) \;\ge\; \vthresh \cdot \vtotal^{\!\star}(\hgt_0),
\qquad \vthresh = 0.25 ,
\]
and $\hgt_c$ denotes the least such height, the \emph{crossing height}. If no such $\hgt_c$ exists,
the report \emph{expires} at $\hgt_0 + \Delta_{\mathrm{vote}}$: it is removed from the open set, its
reporter's capacity is released, and $\mathsf{app\_hash}$ becomes immediately reportable again by
any identity with free capacity. There is no cooldown.
\end{definition}

\begin{definition}[Proposal capacity, the anti-spam rule, \planned]
\label{def:23:capacity}
For $i \in I(\hgt)$ define the capacity $X_i : \mathbb{N} \to \mathbb{N}$,
$X_i(\hgt) \defeq \lvert\{\,j \in N(\hgt) : \mathrm{own}_j = i\,\}\rvert$, the number of confirmed
nodes the identity operates, counted without regard to tier. A report transaction from $i$ is valid
only if the number of $i$'s reports that are open at $\hgt_0$ --- neither expired nor banned --- is
strictly less than $X_i(\hgt_0)$.
\end{definition}

This rule is the whole of the anti-spam design, and it deserves precision rather than praise.
It introduces no new economic primitive: there is no bond, no custody script, no adjudicator that
must learn the outcome to release funds, and therefore none of the open parameters --- size,
custody, destination of a forfeiture --- that a bond drags in. It prices reporting in exactly the
resource that already prices voting, and it inherits the quorum's Sybil resistance without a
separate argument, as follows.

\begin{proposition}[The capacity rule is Sybil-neutral, \planned]
\label{prop:23:capacity}
Let a party $p$ control the node set $N_p \subseteq N(\hgt)$ and partition it arbitrarily among the
owner keys it controls, writing $I_p \defeq \{\mathrm{own}_j : j \in N_p\} \subseteq I(\hgt)$ for
that key set. Then $p$'s total concurrent proposal capacity is
$\sum_{i \in I_p} X_i(\hgt) = \lvert N_p \rvert$, independent of the partition. Increasing it
requires confirming an additional node, hence locking at least $\coll{\tC} = \num{1000}\flux$.
\end{proposition}

\begin{proof}
Each confirmed node has exactly one owner key, so the sets $\{j : \mathrm{own}_j = i\}$ for
$i \in I_p$ partition $N_p$ and their cardinalities sum to $\lvert N_p \rvert$. Re-keying or
re-addressing moves nodes between blocks of the partition and changes no cardinality sum. A node
enters $N(\hgt)$ only through a start transaction against a matured collateral UTXO of a valid tier
amount \srccode{fluxd/\allowbreak{}src/\allowbreak{}coins.cpp}{630--686}, the smallest of which is $\coll{\tC}$.
\end{proof}

\begin{corollary}[Network capacity and sustained report rate, \planned]
\label{cor:23:capacity}
$\sum_{i \in I(\hgt)} X_i(\hgt) = \lvert N(\hgt) \rvert$. At the measured height the network admits
at most \num{6489} concurrently open reports, and --- because capacity is released on expiry and a
report occupies its slot for at most $\Delta_{\mathrm{vote}}$ blocks --- at most \num{6489} reports
per \SI{24}{\hour} in the steady state.
\end{corollary}

The corollary that matters for this section's argument is the unflattering one. Capacity is
distributed as \emph{node counts}, so it is concentrated in the same way and for the same reason
voting power is: the largest node key operates \num{480} of \num{6489} nodes
(\SI{7.40}{\percent} of all capacity; \num{479} under the owner-identity model), and the Gini of node counts is \num{0.716} against
\num{0.839} for weight \srcmeasured{api.runonflux.io/\allowbreak{}daemon/\allowbreak{}listzelnodes}{2026-09-03}. The rule
bounds spam by collateral, which is right; it does not distribute the ability to report any more
evenly than the collateral is distributed, and the concentration measurement of
\Cref{sec:23:measurement} already quantifies exactly how unevenly that is.

It also fixes the price of sustained nuisance, which is worth computing because the number is small.
Keeping one application under a permanently renewed report costs one node --- \num{1000}\flux\ of
\emph{locked, fully recoverable} collateral, not a fee. Keeping every one of the \num{1195}
applications in the deployed registry
\srcmeasured{api.runonflux.io/\allowbreak{}apps/\allowbreak{}globalappsspecifications}{2026-09-03} under permanent report
therefore costs \num{1195000}\flux\ of locked collateral, which is \SI{5.4}{\percent} of the
\num{22128625}\flux\ the eviction threshold itself requires. What that buys is not eviction ---
abstention is ``no'' --- but operator attention and voter fatigue, and the second of those degrades
the liveness the mechanism already has least of.

\begin{theorem}[Public verifiability: nothing is tallied off-chain, \planned]
\label{thm:23:verifiable}
There is a function $\mathrm{Ban}$ from finite chain prefixes to subsets of $\{0,1\}^{256}$,
computable by any node from block data alone, such that under
\Cref{def:23:report,def:23:vote,def:23:decision,def:23:ban} an honest node's enforcement set at
height $\hgt$ equals $\mathrm{Ban}$ applied to its canonical chain at $\hgt - \Delta_{\mathrm{exec}}$.
In particular, two honest nodes holding the same canonical chain hold the same enforcement set
exactly, and no message outside the chain is an input.
\end{theorem}

\begin{proof}
Each of the four objects is a field of a block. Reports and votes are transactions, so
$V_r(\hgt)$ is the set of signers of the valid $\mathsf{MOD\_VOTE}$ transactions in blocks
$[\hgt_0, \hgt]$ and is a function of the prefix. $N(\cdot)$ is consensus state, so
$\vweight{\cdot}(\hgt_0)$ and $\vtotal^{\!\star}(\hgt_0)$ are functions of the prefix; hence so is
$Y_r$, being a finite sum of them, and so is the predicate $Y_r(\hgt_c) \ge
\vthresh\vtotal^{\!\star}(\hgt_0)$. Ban memos are coinbase fields, and their validity condition
(\Cref{def:23:ban}) is that predicate. $\mathrm{Ban}$ is therefore the set of $\mathsf{app\_hash}$
values carried by valid ban memos in the prefix, and enforcement is its image after the
confirmation delay.
\end{proof}

This is the property the section's own verifiability requirement asked for and the property the
message-layer construction could not supply. There, a node holding the chain but not the gossiped
certificate body could not evaluate the predicate at all, so determinism held only ``up to gossip
delay'' and a partitioned node's enforcement set was a function of what it had received. Here the
qualifier disappears. It is also the property the deployed channel of
\Cref{def:23:policy-channel} lacks completely: there the record available to a node for verifying
that a removal was authorised is $\varnothing$.

\subsubsection{Enforcement}
\label{sec:23:enforcement}

\begin{definition}[Ban memo, \planned]
\label{def:23:ban}
A \emph{ban memo} is a field $\mathsf{BAN}\,\|\,H(r)\,\|\,\mathsf{app\_hash}$ carried in the
coinbase transaction of a \PoN{} block at height $\hgt_b$. A block is valid only if it carries
\textbf{at most one} ban memo, and a ban memo is valid only if the report $H(r)$ was carried at
some crossing height $\hgt_c \le \hgt_b - 1$ within its own window, and no valid ban memo for
$H(r)$ appears at any earlier height. The \emph{first eligible height} for a carried report is
therefore $\hgt_c + 1$. Inclusion is \textbf{consensus-required, not discretionary}: every block at
a height at which at least one carried report remains unbanned is invalid unless it carries the
memo for the first such report in the settled order --- least crossing height, then least $H(r)$
\srcintent{evidence/design/08-answers-round-2.md, round 5 and round 18 answers}
(\Cref{prop:23:inclusion-required}). A valid ban memo adds $\mathsf{app\_hash}$ to the ban set,
closes the report, and releases the reporter's capacity.
\end{definition}

\begin{property}[Bounded enforcement throughput and bounded coinbase, \planned]
\label{prop:23:throughput}
At most one application is banned per block, hence at most $\num{2880}$ per \SI{24}{\hour} at
\SI{30}{\second} spacing, and the coinbase grows by at most one fixed-size memo per block
regardless of how many reports are open. This is the same constraint \Cref{sec:07-pnr} respects
when it routes the application-revenue split through the coinbase outputs that already exist rather
than adding new ones.
\end{property}

\begin{property}[The ban set is monotone, \planned]
\label{prop:23:monotone}
No transition specified above removes an entry from the ban set. It is non-decreasing in height,
growing by at most one \SI{32}{\byte} entry per block, and a ban is a chain fact that survives
every future block. \Cref{op:23:reinstatement} is the consequence.
\end{property}

\begin{definition}[Enforcement transition, \planned]
\label{def:23:enforcement}
At every height $\hgt \ge \hgt_b + \Delta_{\mathrm{exec}}$, every \FluxOS{} node that observes a
valid ban memo would: (i) resolve $\mathsf{app\_hash}$ against its global application registry;
(ii) mark the application \texttt{BANNED} and issue the cancel message; (iii) remove any local
instance, reusing the removal path that already exists for the blocklist sweep
\srccode{flux/\allowbreak{}ZelBack/\allowbreak{}src/\allowbreak{}services/\allowbreak{}appSecurity/\allowbreak{}imageManager.js}{644--676}; (iv) refuse to install,
redeploy or soft-redeploy it; (v) stop broadcasting its location. The deployer's prepaid balance
would not be refunded. A memo whose hash resolves to nothing is a no-op.
\end{definition}

Note that the non-refund clause is presently vacuous under deployed code — there is no refund
path at all, and an underpaid registration is silently dropped \srcdoc{\_fluxos.md~\S3.3} — so the
clause acquires content only if the credits and renewal work gives refunds a meaning.

\begin{property}[Ban-memo inclusion is consensus-required, \planned]
\label{prop:23:inclusion-required}
This closes what an earlier draft of this mechanism carried as the single unresolved question about
the enforcement path. The team's decision is that inclusion is obligatory rather than left to
producer discretion: at every height at which a carried report remains unbanned, a producer must
carry the ban memo for the first such report in the settled order, and \Cref{def:23:ban} states the
corresponding validity rule
\srcintent{evidence/design/08-answers-round-2.md, round 5 and round 18 answers}. What it buys is stated plainly:
enforcement latency after the crossing height is bounded by one block for a report with no other
carried report queued ahead of it, not by how long a reluctant or colluding set of producers is
willing to withhold it. Which \emph{one} report's memo a block
must carry, when more than one crosses in the same slot, is the ordering question
immediately below, and the obligation binds over that order.
\end{property}

\settlednote{Least crossing height, then least $H(r)$; the memo is carried in the coinbase (\Cref{tab:23-recommended}). Ordering when several reports are carried at the same height, and the carriage of the memo
in the coinbase. Domains: ordering $\in \{$least $H(r)$; least $\hgt_0$; least crossing height then
least $H(r)$; producer's choice$\}$; carriage $\in \{$a zero-value \texttt{OP\_RETURN} output; a
field in the coinbase \texttt{scriptSig}$\}$. Trade-off: with one memo per block a queue forms
whenever two reports cross in the same slot, and the settled obligatory-inclusion rule cannot be
enforced without a deterministic order over that queue; on carriage, an extra coinbase \emph{output}
interacts with the two independent checks that already scan coinbase outputs --- the dev-fund
minimum \srccode{fluxd/\allowbreak{}src/\allowbreak{}main.cpp}{1231--1249} and the tier-payout check
\srccode{fluxd/\allowbreak{}src/\allowbreak{}fluxnode/\allowbreak{}fluxnode.cpp}{837--875} --- while a \texttt{scriptSig} field does not.}

The lower bound on the penalty follows from the deployed pricing table. The forfeited amount is the
unused fraction of the deployment price, floored at \texttt{minPrice} $=$ \num{0.01}\flux\ per
month-equivalent and pro-rated to the validator's height-gated minimum term --- \num{5000} blocks,
then \num{100} from height \num{1630040}, then \num{1} from height \num{1964447}
\srccode{flux/\allowbreak{}ZelBack/\allowbreak{}src/\allowbreak{}services/\allowbreak{}appRequirements/\allowbreak{}appValidator.js}{852--862},
while the live endpoint still reports \num{5000}
\srcmeasured{api.runonflux.io/\allowbreak{}apps/\allowbreak{}deploymentinformation}{2026-09-03} (conflict C33) --- against a
post-\PoN{} default of \num{88000}, then re-ceiled to two decimals. A minimally provisioned
application would therefore forfeit \num{0.01}\flux\ under either minimum.

\begin{property}[Delay is fixed at one slot under the settled rule, \planned]
\label{prop:23:inclusion}
Under \Cref{prop:23:inclusion-required}, enforcement latency after a report crosses threshold does
not depend on which node produces the next block: the block at the first eligible height is invalid
without the memo, so $\delta = 1$ whenever no other carried report is queued ahead of it (one memo
per block, \Cref{def:23:ban}). This replaces what a discretionary rule would have
cost. Assuming \PoN{} slot eligibility is uniform over confirmed nodes and tier-independent
\srccode{fluxd/\allowbreak{}src/\allowbreak{}pon/\allowbreak{}pon.cpp}{400--433} and that successive slots' producers are drawn
independently, a discretionary rule under which a fraction $f$ of confirmed nodes decline to include
a due memo would have produced a geometric delay --- expected $1/(1-f)$ slots,
$\Pr[\text{delay} > k] = f^{\,k}$ --- diverging only as $f \to 1$. At the measured distribution's
largest identity, $f = 0.074$, that would have been \num{1.08} slots in expectation, about
\SI{32}{\second}; at $f = 0.5$, two slots. The counterfactual shows the discretionary rule would not
have been a practical censorship vector at the measured distribution either, which makes the settled
rule the more conservative choice rather than a strictly necessary one --- it removes the dependence
on $f$ entirely instead of merely making it small.
\end{property}

The vote-transaction censorship question is separate and is unaffected by
\Cref{prop:23:inclusion-required}: a vote needs only one non-declining producer across the whole
\num{2880}-slot window to be included, so under the same independence assumption a coalition holding
fraction $f$ of nodes excludes a given vote from the entire window with probability
$f^{\num{2880}}$, negligible for every $f < 1$. That the \emph{window}, not any single slot, is what
protects a vote is a different, and looser, guarantee than the exact one-slot bound the ban memo now
has by consensus rule.

One constraint this subsection exports to consensus, in the same class as the constraint
\Cref{sec:07-pnr} carries on the fee pool. The 2-of-4 emergency key set can produce a block outside
normal \PoN{} eligibility \srccode{fluxd/\allowbreak{}src/\allowbreak{}emergencyblock.cpp}{24--30,183--190}, and that block carries a
coinbase. \textbf{The ban-memo validity rule of \Cref{def:23:ban} must therefore bind on the
emergency path as well}, or those keys acquire the ability to ban an application with no report,
no votes and no threshold --- a capability materially larger than the one the emergency path exists
to provide.

\begin{remark}[Emergency blocks are subject to the coinbase rules --- verified]
\label{rem:emergency-coinbase}
This was carried as an open verification, because a ban memo carried in the coinbase is only as
sound as the weakest block-production path, and the emergency path bypasses eligibility entirely
(\S\ref{sec:consensus}). It is now checked against source and the answer is favourable:
\texttt{Is\allowbreak{}Emergency\allowbreak{}Block} is consulted in six places in consensus code, all of them header
validation --- to permit a larger block signature and to skip the proof-of-node check in
\texttt{Check\allowbreak{}Block\allowbreak{}Header} \srccode{fluxd/src/main.cpp}{5333,5339}, in the equivalent header path
\srccode{fluxd/src/main.cpp}{2778}, when reloading headers from disk
\srccode{fluxd/src/txdb.cpp}{557}, and in the \PoN{} header checks that skip the proof of node and
validate the emergency signatures and activation gate \srccode{fluxd/src/pon/pon.cpp}{205,236--250}.
It appears \emph{nowhere} in \texttt{ConnectBlock},
\texttt{Contextual\allowbreak{}Check\allowbreak{}Block} or \texttt{Contextual\allowbreak{}Check\allowbreak{}Transaction} \shipped.

\textbf{So the emergency bypass is narrow and does not reach the coinbase.} An emergency block skips
the proof of node and the ordinary signature-size limit, and is otherwise validated by every rule a
normal block is: BIP-30 duplicate-transaction rejection
\srccode{fluxd/src/main.cpp}{3868--3875}, the deterministic FluxNode payout check
\srccode{fluxd/src/main.cpp}{3877--3880}, and the full contextual transaction checks. A ban memo
specified as a coinbase constraint is therefore enforced on emergency blocks too, and the constraint
in \Cref{sec:23:mechanism} can be stated as a claim about deployed behaviour rather than as a
requirement awaiting confirmation.
\end{remark}

\subsubsection{Placement, and what fee-lessness costs}

\begin{table}[H]
\centering
\small
\begin{tabular}{@{}p{3.1cm}p{4.0cm}p{7.4cm}@{}}
\toprule
object & location & reason \\
\midrule
report & \textbf{chain}, indexed & canonical $\hgt_0$; freezes the electorate; discoverable by any node \\
votes & \textbf{chain}, fee-less & one transaction per identity per report; confirmation height is the canonical clock \\
tally $Y_r$ & \textbf{chain-derived}, in \fluxd{} & a pure function of the block sequence; recomputed, never journaled \\
decision & \textbf{chain-derived} & the predicate of \Cref{def:23:decision}, evaluated by every validator \\
ban & \textbf{chain}, coinbase memo & one per block; canonical $\hgt_b$, hence a canonical enforcement height \\
enforcement & \FluxOS{}, at $\hgt_b + \Delta_{\mathrm{exec}}$ & where removal already happens \\
\bottomrule
\end{tabular}
\caption{Placement of each object in the specified mechanism (\planned). Nothing in the decision
path leaves the chain, which is the difference between this design and the message-layer
construction it replaces.}
\end{table}

An earlier draft of this mechanism rejected on-chain voting for two reasons: it would cost a
fee-bearing transaction per voter, and it would require every voting key to hold spendable \flux,
which is a key-hygiene hazard the network does not currently have. \textbf{The fee-less rule
retires both objections.} An operator with no spendable balance at the key that carries their
collateral can still vote, so no operator is disenfranchised by their own treasury arrangements ---
which is not a small point for a mechanism whose safety argument already depends on the dispersed
majority being able to participate at all.

\begin{remark}[Fee-lessness is structural, not a premise --- verified]
\label{rem:feeless-verified}
The fee-less rule was carried as resting on the team's stated premise that \FluxNode{} start and
confirm transactions pay no fee. It is stronger than a premise: it is enforced by consensus.
\texttt{Check\allowbreak{}Transaction} rejects any \FluxNode{} transaction whose \texttt{vin}, \texttt{vout},
\texttt{vJoinSplit}, \texttt{v\allowbreak{}Shielded\allowbreak{}Output} or \texttt{vShieldedSpend} is non-empty, with reject
reason \texttt{bad-\allowbreak{}txns-\allowbreak{}fluxnode-\allowbreak{}tx-\allowbreak{}not-\allowbreak{}empty} \shipped
\srccode{fluxd/src/main.cpp}{1751--1755}. A transaction with no inputs and no outputs
\emph{cannot} pay a fee, since a fee is the difference between them.

So a vote following the start-transaction shape (\Cref{tab:23-recommended}) inherits fee-lessness by
construction rather than by policy, and cannot be made fee-bearing without changing the transaction
shape itself. That is why the relay policy needs a non-fee scarcity term --- the per-identity mempool
quota recommended in \Cref{tab:23-recommended} --- and why no fee-based anti-spam mechanism is
available here even in principle.
\end{remark}

What fee-lessness costs is the price signal that would otherwise bound volume, so the bound must
come from the capacity rule instead, and it is worth seeing the worst case. With every capacity
slot used --- \num{6489} concurrently open reports by \Cref{cor:23:capacity} --- and every one of
the \num{869} owner identities voting on every report, the window carries
$\num{6489} \times \num{869} = \num{5638941}$ vote transactions, about \num{1958} per block. At an
assumed \SI{250}{\byte} per vote transaction that is \SI{0.49}{\mega\byte} per block against a
\SI{2}{\mega\byte} block-size limit \srccode{fluxd/\allowbreak{}src/\allowbreak{}consensus/\allowbreak{}consensus.h}{23}, or
\SI{1.41}{\giga\byte} of the \SI{5.76}{\giga\byte} a day of blocks can carry --- roughly a quarter
of the chain's entire capacity, contributed by transactions that pay nothing \srcinferred. The
transaction format is not specified, so the \SI{250}{\byte} figure is an assumption and the
conclusion is a shape rather than a number: \emph{the concurrent-proposal cap is load-bearing for
block space, not merely for operator attention}, and it should be chosen with that in mind.

\subsubsection{State machine and node-side procedure}

\begin{table}[H]
\centering
\small
\begin{tabular}{@{}p{2.0cm}p{5.6cm}p{1.9cm}p{4.6cm}@{}}
\toprule
from & guard & to & effect \\
\midrule
\textsc{submitted} & report tx confirms at $\hgt_0$; $\sigma_r$ verifies;
  $\mathrm{own}_r \in I(\hgt_0)$; open reports of $\mathrm{own}_r < X_{\mathrm{own}_r}(\hgt_0)$;
  no \textsc{open} report on $\mathsf{app\_hash}$; $\mathsf{app\_hash} \notin \mathcal{B}$
  & \textsc{open} & electorate frozen at $\hgt_0$; one capacity slot consumed \\
\textsc{submitted} & any guard fails & \textsc{invalid} & the transaction is rejected; nothing
  is spent, since there is no fee and no bond \\
\textsc{open} & $Y_r(\hgt_c) \ge \vthresh\vtotal^{\!\star}(\hgt_0)$, $\hgt_c \le \hgt_0 +
  \Delta_{\mathrm{vote}}$ & \textsc{carried} & first eligible ban height is $\hgt_c + 1$ \\
\textsc{open} & $\hgt > \hgt_0 + \Delta_{\mathrm{vote}}$, threshold not reached & \textsc{expired}
  & capacity released; $\mathsf{app\_hash}$ immediately re-reportable; no cooldown \\
\textsc{carried} & the first block at $\hgt_b \ge \hgt_c + 1$ at which the report is first in the
  settled order carries its ban memo, consensus-required (at most one per block) & \textsc{banned}
  & $\mathsf{app\_hash}$ enters the ban set; capacity released \\
\textsc{banned} & $\hgt \ge \hgt_b + \Delta_{\mathrm{exec}}$ & \textsc{enforced} &
  \Cref{def:23:enforcement} \\
\textsc{enforced} & --- & --- & no transition exists (\Cref{op:23:reinstatement}) \\
\bottomrule
\end{tabular}
\caption{Report lifecycle (\planned). Every guard is a predicate over the canonical chain, so every
transition is evaluated identically by every validator; there is no state that a node can be
missing.}
\end{table}

\begin{algorithm}[H]
\caption{Consensus-side validation and tally, per block (\planned, \fluxd)}
\KwIn{block $B$ at height $\hgt$; the confirmed node set $N(\cdot)$ and its per-block undo data
\srccode{fluxd/\allowbreak{}src/\allowbreak{}fluxnode/\allowbreak{}fluxnodecachedb.cpp}{14--17}; the open-report table $\mathcal{O}$ and
ban set $\mathcal{B}$ as recomputed from the chain prefix}
\KwOut{accept or reject $B$; updated $\mathcal{O}$ and $\mathcal{B}$}
\tcc{Invariant: $\mathcal{O}$ and $\mathcal{B}$ are pure functions of the chain prefix ending at
$B$. Nothing is written that a resyncing node could not recompute, and no input comes from outside
the chain — which is what makes a separate moderation database unnecessary, and, per the fee-pool
recommendation of \Cref{sec:07-pnr}, undesirable.}
\tcc{Failure mode: consensus validates form and authorisation, never the referent — a report naming
a hash with no application behind it is valid here and a no-op in \FluxOS. Fee-less transactions
carry no price signal, so mempool admission must be bounded by capacity, not by fee. On a reorg,
$\mathcal{O}$ and $\mathcal{B}$ are recomputed with the chain and need no repair path.}
\ForEach{transaction $\mathrm{tx} \in B$}{
  \If{$\mathrm{tx}$ is $\mathsf{MOD\_REPORT}$}{
    \lIf{$\sigma_r$ invalid \textbf{or} $\mathrm{own}_r \notin I(\hgt)$ \textbf{or}
      $\lvert\{\text{open reports of } \mathrm{own}_r\}\rvert \ge X_{\mathrm{own}_r}(\hgt)$
      \textbf{or} an open report exists on $\mathsf{app\_hash}$
      \textbf{or} $\mathsf{app\_hash} \in \mathcal{B}$}{\textbf{reject} $B$}
    insert $H(r)$ into $\mathcal{O}$ with $\hgt_0 \leftarrow \hgt$;\quad
      $Y_r \leftarrow 0$;\quad $V_r \leftarrow \varnothing$\;
  }
  \If{$\mathrm{tx}$ is $\mathsf{MOD\_VOTE}$ on $H(r)$ by $i$}{
    \lIf{$H(r) \notin \mathcal{O}$ \textbf{or} $\sigma_v$ invalid \textbf{or} $i \notin I(\hgt_0)$
      \textbf{or} $\hgt > \hgt_0 + \Delta_{\mathrm{vote}}$ \textbf{or} $i \in V_r$}{\textbf{reject} $B$}
    $V_r \leftarrow V_r \cup \{i\}$;\quad
      \lIf{$\mathsf{decision} = \mathsf{evict}$}{$Y_r \leftarrow Y_r + \vweight{i}(\hgt_0)$}
  }
}
\ForEach{$H(r) \in \mathcal{O}$ with $\hgt > \hgt_0 + \Delta_{\mathrm{vote}}$ and
  $Y_r < \vthresh\,\vtotal^{\!\star}(\hgt_0)$}{
  delete $H(r)$ from $\mathcal{O}$ (\textsc{expired}); release one capacity slot of
    $\mathrm{own}_r$\;
}
\lIf{some $H(r) \in \mathcal{O}$ is carried with crossing height ${} \le \hgt - 1$ \textbf{and} the
  coinbase of $B$ carries no ban memo for the first such report in the settled order (least crossing
  height, then least $H(r)$)}{\textbf{reject} $B$}
\If{the coinbase of $B$ carries a ban memo for $H(r)$}{
  \lIf{$B$ carries more than one memo \textbf{or} $H(r) \notin \mathcal{O}$ \textbf{or}
    $Y_r < \vthresh\,\vtotal^{\!\star}(\hgt_0)$ \textbf{or} the crossing height is not
    ${} \le \hgt - 1$ \textbf{or} $\mathsf{app\_hash} \in \mathcal{B}$}{\textbf{reject} $B$}
  $\mathcal{B} \leftarrow \mathcal{B} \cup \{\mathsf{app\_hash}\}$; delete $H(r)$ from
    $\mathcal{O}$; release the capacity slot\;
}
\end{algorithm}

\begin{algorithm}[H]
\caption{Ban observation and purge, per node (\planned, \FluxOS)}
\KwIn{canonical chain; confirmation delay $\Delta_{\mathrm{exec}}$; local application set
$A_{\mathrm{loc}}$; the global specification registry}
\KwOut{registry state and $A_{\mathrm{loc}}$ for every banned $\mathsf{app\_hash}$}
\tcc{Invariant (exact determinism): enforcement is a function of the canonical chain alone, so two
honest nodes on the same chain reach the same enforcement set — not merely up to gossip delay, as
in the message-layer construction this replaces.}
\tcc{Failure mode (fail-safe, with one asymmetry): a memo whose hash resolves to no registry entry
is a no-op, and a node that has not yet reached $\hgt_b + \Delta_{\mathrm{exec}}$ leaves the
application up. The asymmetry is that purging is not locally reversible: a reorg deeper than
$\Delta_{\mathrm{exec}}$ that removes the ban block leaves an application deleted that the chain no
longer bans, and it must be redeployed. This is why $\Delta_{\mathrm{exec}}$ is a reorg parameter
and not a policy one.}
\ForEach{valid ban memo observed at height $\hgt_b$}{
  \lIf{$\hgt < \hgt_b + \Delta_{\mathrm{exec}}$}{\textbf{wait}}
  \lIf{$\mathsf{app\_hash}$ resolves to no registry entry}{record and \textbf{continue} (no-op)}
  mark the application \texttt{BANNED}; issue the cancel message\;
  \lIf{an instance is installed locally}{remove it by the deployed removal path; refuse subsequent
    install, redeploy and soft-redeploy; stop location broadcasts}
}
\end{algorithm}

The voting window is settled at $\Delta_{\mathrm{vote}} = \num{2880}$ blocks
\srcintent{evidence/design/08-answers-round-2.md}, and the reason given is worth recording because
it cuts against turnout: removal often needs to be fast, and a \SI{24}{\hour} window is the shortest
one that still lets an operator who checks daily participate. Everything \Cref{sec:23:safety} says
about liveness is said against that number.

\settlednote{$\Delta_{\mathrm{exec}} = \num{120}$ blocks (\Cref{tab:23-recommended}). $\Delta_{\mathrm{exec}}$, the delay between the ban block $\hgt_b$ and enforcement.
Domain: $[0,\ \num{2880}]$ blocks. Trade-off: it must exceed plausible reorg depth --- \PoN{} fork
choice is a block-count race because every block carries fixed $2^{40}$ chainwork
\srccode{fluxd/\allowbreak{}src/\allowbreak{}pow.cpp}{284--290}, and the deployed reorg bound outside the activation window is
40 blocks \srccode{fluxd/\allowbreak{}src/\allowbreak{}main.h}{74} --- because a purge is not locally reversible, while every
block of delay is a block of harm uptime. It is also the mechanism's only post-decision safety
margin, the analogue of a governance timelock (\Cref{sec:23:related}). Note that
$\Delta_{\mathrm{exec}} = 0$ is a defensible choice only if a purged application returning after a
reorg is acceptable.}

\begin{assumption}[Model assumptions for the specified mechanism]
\label{ass:23:mechanism}
As fixed in the system-model section, and specialised here. (A1) \fluxd's deterministic node set at
any height is consensus state and is agreed by all honest nodes. (A2) Owner keys are controlled by
the parties the mechanism intends to enfranchise; the consequence of this failing is treated in
\Cref{sec:23:measurement}. (A3) A transaction broadcast during the window is included by some
producer before the window closes --- which \Cref{prop:23:inclusion} shows requires only that not
every producer censors it, over \num{2880} slots. (A4) A chain reorg deeper than
$\Delta_{\mathrm{exec}}$ does not occur. (A5) Signature forgery is infeasible. No assumption is made
about turnout; that is the point of \Cref{thm:23:turnout}. Note what is \emph{not} assumed, and was
assumed by the construction this replaces: no property of the \FluxOS{} gossip mesh is needed
anywhere in the decision path.
\end{assumption}

The mechanism's invariants follow directly and are stated without separate proof: $Y_r$ is
non-negative and non-decreasing within the window; weights and the reference total are functions of
$\hgt_0$ alone; one vote per identity per report, enforced by chain order; a vote verifies against
exactly one $(H(r), \hgt_0)$; each report resolves exactly once, into \textsc{expired} or
\textsc{banned}, and releases exactly one capacity slot when it does; a node that cannot resolve a
ban hash does not remove anything; the ban set is monotone (\Cref{prop:23:monotone}); the coinbase
grows by at most one memo per block; and no transition reads or writes a collateral UTXO, the node
set, or a node's payment rank. That last one is the settled promise that node collateral is never
slashed, and it is what removes the most obvious griefing concern — and, as \Cref{sec:23:dilution}
shows, what makes voting close to costless. One invariant of the earlier draft does \emph{not}
survive unconditionally: that an update during the window is neither an escape nor a second claim
holds only under the identity reading of $\mathsf{app\_hash}$ --- which is now the settled one,
$(\mathsf{app\_name}, \mathsf{owner})$ (\Cref{tab:23-recommended}), so the invariant does hold on the
settled configuration. It is stated conditionally because it is a property of that choice rather than
of the mechanism.

\begin{property}[No reinstatement: a ban is final, \planned]
\label{op:23:reinstatement}
The specified design contains no transition that removes an application from the ban set and no
transition that returns a forfeited prepaid balance, and this is a settled position rather than a
gap: the team's decision is that a ban is final, and the remedy is redeployment as a \textbf{new}
application \srcintent{evidence/design/08-answers-round-2.md, round 5 answers}. Both absences are
irreversible by construction, which is what makes the decision coherent rather than merely
convenient: the forfeiture is a non-event for a pool that was credited at payment time
(\Cref{prop:23:neutral}), and the ban is a chain fact carried in a coinbase that no later block
retracts (\Cref{prop:23:monotone}). If a quorum is wrong --- honestly wrong on mistaken evidence, or
dishonestly wrong under \Cref{thm:23:collusion} --- redeployment under a new name is the entire
remedy, and it preserves none of the old instance's reachability, since \FDM{} and \PDNS{} names are
keyed per application name. A construction that reversed a ban would in any case have to satisfy the
requirements an earlier draft of this section raised --- a chain object, so
\Cref{thm:23:verifiable} is not lost for the reversal direction; not a symmetric quorum over the
same electorate, or the coalition that can ban can also block reinstatement
(\Cref{sec:23:remedies}); and an interaction with the forfeiture, since reinstating an application
whose term was consumed reinstates it into nothing --- which is a substantial part of why the
team's simpler position is defensible rather than merely a shortcut.

One residual is worth stating precisely rather than absorbing into the decision silently. If
$\mathsf{app\_hash}$ resolves to the application's \emph{content} --- a specification or image hash
--- rather than to its $(\mathsf{app\_name}, \mathsf{owner})$ \emph{identity}, a redeployment that is
otherwise identical carries the same hash and is therefore immediately re-bannable: the hash is
already in the ban set, with no fresh report and no fresh vote required. Whether that is the
intended reading or an accident of the keying choice is exactly the enforcement-scope question left
open earlier in this section (``What $\mathsf{app\_hash}$ hashes, which is also the enforcement
scope,'' \Cref{sec:23:mechanism}); this paper cross-references it here rather than deciding it
twice, since what a ban targets and what redeployment resets are the same question asked from
opposite ends.
\end{property}

% =====================================================================
\subsection{Safety and liveness, separated}
\label{sec:23:safety}

\begin{definition}[Safety and liveness for this mechanism]
\label{def:23:safety-liveness}
\emph{Safety} (against wrongful eviction): an application is not banned unless a set of identities
of combined weight at least $\vthresh\vtotal^{\!\star}(\hgt_0)$ has confirmed $\mathsf{evict}$ votes
on one report. \emph{Liveness}: an application that a set of identities of combined weight at least
$\vthresh\vtotal^{\!\star}$ wishes to remove is removed within
$\Delta_{\mathrm{vote}} + \Delta_{\mathrm{exec}} + \delta$ blocks, where $\delta$ is the number of
slots between the crossing height and the block that carries the memo. Under the settled,
consensus-required inclusion rule $\delta$ is one plus the number of unbanned carried reports ahead
of it in the settled order, so $\delta = 1$ whenever no other carried report is queued ahead of
it (\Cref{prop:23:inclusion-required,prop:23:inclusion}).
\end{definition}

\begin{lemma}[Turnout bound]
\label{lem:23:turnout}
Let $P_r \subseteq I(\hgt_0)$ be the set of identities whose vote on $r$ confirms within the window.
Then $Y_r(\hgt) \le \sum_{i \in P_r} \vweight{i}(\hgt_0)$ for all $\hgt$ in the window, and if
$\sum_{i \in P_r} \vweight{i}(\hgt_0) < \vthresh\,\vtotal^{\!\star}(\hgt_0)$ then the report cannot
be carried and expires.
\end{lemma}

\begin{proof}
$V_r(\hgt) \subseteq P_r$ by construction and all weights are non-negative, so
$Y_r(\hgt) = \sum_{i \in V_r(\hgt)} \vweight{i} \le \sum_{i \in P_r} \vweight{i}$. Carrying requires
$Y_r(\hgt_c) \ge \vthresh\vtotal^{\!\star}$ for some $\hgt_c$ in the window, which the displayed
inequality forbids; \Cref{def:23:decision} then expires the report.
\end{proof}

\begin{theorem}[Turnout is a liveness parameter and never a safety parameter]
\label{thm:23:turnout}
Under \Cref{ass:23:mechanism}, for any report $r$ and any participation set $P_r$: adding a
non-voter to the electorate can only increase $\vtotal^{\!\star}(\hgt_0)$ and hence the weight
required for a certificate; removing a voter can only decrease $Y_r$. Consequently low turnout
makes eviction strictly harder, never easier, and no low-turnout election can produce an eviction
that a full-turnout election with the same $\mathsf{evict}$ votes would not also produce.
\end{theorem}

\begin{proof}
The decision predicate is $Y_r(\hgt_c) \ge \vthresh\,\vtotal^{\!\star}(\hgt_0)$. The right-hand side
depends on the electorate at $\hgt_0$ and not at all on who votes; adding an identity $i'$ to
$I(\hgt_0)$ increases $\vtotal(\hgt_0)$ by $\vweight{i'} \ge 0$ and hence weakly increases
$\vtotal^{\!\star}$ and the required threshold. The left-hand side is a sum over the identities that
actually voted $\mathsf{evict}$; deleting one removes a non-negative term. Therefore the predicate
is monotone decreasing in the electorate and monotone increasing in the $\mathsf{evict}$ set, and
the two are independent. Fixing the $\mathsf{evict}$ set and varying turnout among the remaining
identities changes neither side.
\end{proof}

The contrast is the point. Had the threshold been defined over votes cast — that is,
$Y_r / \sum_{i \in P_r} \vweight{i} \ge \vthresh$ — then a single identity of any weight could evict
any application by being the only voter, and low turnout would convert from a liveness limitation
into a safety hole. The settled rule excludes this. Stated in the design's own terms: the
\SI{25}{\percent} figure is a liveness/safety dial, abstention counts against removal, and the
design deliberately errs toward leaving content up rather than removing it wrongly. That is the
defensible position for a censorship-resistant network, and it is the strongest property this
mechanism has.

It is paid for in liveness, and the bill is legible.

\begin{proposition}[Liveness]
\label{prop:23:liveness}
If identities of combined weight at least $\vthresh\vtotal^{\!\star}(\hgt_0)$ broadcast
$\mathsf{evict}$ votes on $r$ and those transactions confirm by
$\hgt_0 + \Delta_{\mathrm{vote}}$, then the report is carried at some
$\hgt_c \le \hgt_0 + \Delta_{\mathrm{vote}}$, and the application is removed by
$\hgt_c + \delta + \Delta_{\mathrm{exec}}$ with $\delta$ as in
\Cref{def:23:safety-liveness}.
\end{proposition}

\begin{proof}
By hypothesis $Y_r$ reaches $\vthresh\vtotal^{\!\star}(\hgt_0)$ at the height at which the last of
those votes confirms, which is within the window; that height is $\hgt_c$ or later than the least
such height, and \Cref{def:23:decision} takes the least. A producer at some height
$\hgt_b = \hgt_c + \delta$ carries the memo, valid by \Cref{def:23:ban}; by (A4) it is not reorged
out, and by \Cref{thm:23:verifiable} every honest node computes the same ban set and applies
\Cref{def:23:enforcement} at $\hgt_b + \Delta_{\mathrm{exec}}$.
\end{proof}

Note what dropped out of the hypothesis. The message-layer construction needed a party willing to
assemble a certificate and pay for an anchoring transaction --- a coordination step with no assigned
actor and no incentive attached to it. Here there is nothing to assemble: the votes are already in
the chain, and the count is a function every validator already evaluates. That is a genuine
liveness improvement and it comes from the same choice that delivers
\Cref{thm:23:verifiable}.

What that hypothesis costs at the measured distribution:

\begin{table}[H]
\centering
\small
\begin{tabular}{@{}lrr@{}}
\toprule
who abstains & weight remaining & share of the remainder that must vote $\mathsf{evict}$ \\
\midrule
nobody & \SI{100}{\percent} & \SI{25.0}{\percent} \\
the 3 largest keys & \SI{74.9}{\percent} (\num{66314.5} units) & \textbf{\SI{33.4}{\percent}} \\
the 16 largest keys & \SI{49.9}{\percent} (\num{44193.0} units) & \textbf{\SI{50.1}{\percent}} \\
\bottomrule
\end{tabular}
\caption{Turnout required among the remainder to reach $\vthresh = 0.25$, by abstaining prefix,
computed from the measured weight distribution
\srcmeasured{api.runonflux.io/\allowbreak{}daemon/\allowbreak{}listzelnodes}{2026-09-03}.}
\end{table}

The design intake's own premise is that most operators do not vote and will not follow votes
continuously \srcintent{evidence/design/04-moderation.md}. Combined with
\Cref{lem:23:turnout}, the operational reading is that at $\vthresh = 0.25$ liveness depends on the
large keys participating, because the dispersed remainder is unlikely to reach a third of itself.
The same identities that would make the mechanism unsafe are the ones that would make it live.

The \SI{24}{\hour} window makes that bill due in one day rather than seven, and the honest way to
present the consequence is not as a defect.

\begin{proposition}[Expiry is the mechanism working]
\label{prop:23:expiry}
Carrying a report requires identities holding at least
$\vthresh\vtotal^{\!\star}(\hgt_0) = \num{22128.6}$ units at the measured distribution --- the
equivalent of \num{554} Stratus nodes --- to get $\mathsf{evict}$ transactions \emph{confirmed}
inside \num{2880} blocks. Whenever that does not happen the report expires, and by
\Cref{lem:23:turnout} expiry is exactly the event ``no set of the required weight voted''. An
expired report removes nothing, penalises nobody, and returns the capacity slot; the application
stays up.
\end{proposition}

Against the abstention-counts-as-no rule this is a safety property, not a failure. A missed window
is the mechanism declining to act on insufficient agreement, which is the direction the settled dial
points, and the design's own premise about operator participation implies that most reports will
end this way. The paper states the expectation in those terms rather than presenting a
low-success-rate mechanism as a broken one: what a \SI{24}{\hour} window at
$\vthresh = 0.25$ buys is a fast path for the cases where the large holders agree quickly, and
nothing else. The cost of that choice is stated in \Cref{prop:23:liveness}'s hypothesis and in the
table above; the benefit is that a genuinely harmful workload can be gone in a day rather than a
week.

% =====================================================================
\subsection{The measurement}
\label{sec:23:measurement}

Collateral weighting is deliberate and settled, and it is settled for a good reason: voting power
tracks economic stake because collateral is what is at risk. It delivers Sybil resistance as a
property of the construction rather than as a hope.

\begin{property}[Sybil resistance]
\label{prop:23:sybil}
Voting weight cannot be increased without locking additional collateral. Re-keying, re-addressing,
splitting one node's collateral across many node records, or registering additional identities
never changes $\sum_i \vweight{i}$ under one party's control, because $\vweight{i}$ is a sum over
collateral amounts attested by consensus state and each collateral UTXO backs exactly one confirmed
node \srccode{fluxd/\allowbreak{}src/\allowbreak{}fluxnode/\allowbreak{}fluxnode.cpp}{445--453}.
\end{property}

This is the resource-cost answer to the identity-multiplication attack of
\citet{douceur2002sybil}, and it is the reason the weighting is settled.
\Cref{prop:23:capacity} is the identical argument applied to proposal capacity, which is why the
anti-spam rule inherits this property instead of needing one of its own: both quantities are sums
over confirmed nodes, and a confirmed node costs collateral. The same choice is what
concentrates the mechanism, and the concentration is measurable.

\begin{definition}[Collusion set]
\label{def:23:collusion}
For a threshold $\vthresh \in (0,1]$,
\[
\kappa(\vthresh) \;\defeq\; \min\Bigl\{\,\lvert S\rvert \ :\ S \subseteq I(\hgt_0),\
\textstyle\sum_{i \in S} \vweight{i}(\hgt_0) \;\ge\; \vthresh\,\vtotal(\hgt_0)\,\Bigr\}.
\]
Because weights are fixed and non-negative, the minimum is attained by a greedy prefix of
identities in decreasing weight order.
\end{definition}

\begin{theorem}[Safety is exactly a collusion bound]
\label{thm:23:collusion}
Under \Cref{ass:23:mechanism}, an application can be wrongly evicted only if some set of at
least $\kappa(\vthresh)$ identities all sign $\mathsf{evict}$: no smaller set can, at any turnout.
Conversely, the $\kappa(\vthresh)$ heaviest identities can, at any turnout, whenever
$\vtotal^{\!\star}(\hgt_0) = \vtotal(\hgt_0)$.
\end{theorem}

\begin{proof}
($\Rightarrow$) The voter set that carries a report is a set $S$ with $\sum_{i \in S}\vweight{i} \ge
\vthresh\vtotal^{\!\star} \ge \vthresh\vtotal$, so $\lvert S\rvert \ge \kappa(\vthresh)$ by
minimality in \Cref{def:23:collusion}. ($\Leftarrow$) The greedy prefix of size $\kappa(\vthresh)$
has weight at least $\vthresh\vtotal = \vthresh\vtotal^{\!\star}$ under the stated condition, so its
confirmed $\mathsf{evict}$ votes carry the report whatever every other identity does, since by \Cref{lem:23:turnout} other identities can
only add to $Y_r$ and never subtract from it, and the threshold does not depend on them. The claim
``at any turnout'' is \Cref{thm:23:turnout}.
\end{proof}

The mechanism therefore has the safety property ``no coalition smaller than $\kappa(\vthresh)$ can
censor''. Whether that is a property or an assumption about people depends entirely on the number.
Computed from the \num{6489} raw node records:

\begin{table}[H]
\centering
\small
\begin{tabular}{@{}lrrrrr@{}}
\toprule
$\vthresh$ & required weight (units) & $\kappa$ by owner identity (869)
  & $\kappa$ by node key (819) & Stratus nodes alone & of Stratus set \\
\midrule
\SI{10}{\percent}   & \num{8851.5}  & \textbf{1} & \textbf{1} & 222 & \SI{13.6}{\percent} \\
\SI{15}{\percent}   & \num{13277.2} & 2 & 2 & 332 & \SI{20.4}{\percent} \\
\SI{20}{\percent}   & \num{17702.9} & 3 & 2 & 443 & \SI{27.2}{\percent} \\
\textbf{\SI{25}{\percent}} & \textbf{\num{22128.6}} & \textbf{4} & \textbf{3} & \textbf{554}
  & \textbf{\SI{34.1}{\percent}} \\
\SI{30}{\percent}   & \num{26554.4} & 6 & 5 & 664 & \SI{40.8}{\percent} \\
\SI{33.3}{\percent} & \num{29475.3} & 8 & 6 & 737 & \SI{45.3}{\percent} \\
\SI{40}{\percent}   & \num{35405.8} & 12 & 10 & 886 & \SI{54.5}{\percent} \\
\SI{50}{\percent}   & \num{44257.3} & 19 & 16 & \num{1107} & \SI{68.0}{\percent} \\
\SI{51}{\percent}   & \num{45142.4} & 20 & 17 & \num{1129} & \SI{69.4}{\percent} \\
\bottomrule
\end{tabular}
\caption{$\kappa(\vthresh)$ under the settled weighting, by identity proxy, from
\srcmeasured{api.runonflux.io/\allowbreak{}daemon/\allowbreak{}listzelnodes}{2026-09-03} (\num{6489} records, chain tip
\num{2917115}; this table and every other figure this section computes from that snapshot are
reproduced by \texttt{scripts/\allowbreak{}22-moderation-tables.\allowbreak{}py}). The operative column is the first: the vote is signed with the owner key, whose
public proxy is \texttt{payment\_address}. The node-key column is retained only because it is the
lower of the two and therefore the conservative reading; \texttt{pubkey} is the VPS message-signing
key and is not the vote key \srccode{fluxd/\allowbreak{}src/\allowbreak{}fluxnode/\allowbreak{}fluxnode.h}{154}. Merging addresses into
clusters by shared node key or shared collateral funding transaction (797 clusters) gives
$1 / 2 / 3 / 16$ at \SI{10}{\percent} / \SI{15}{\percent} / \SI{25}{\percent} / \SI{50}{\percent};
the intermediate rows were not computed for that model. Nimbus alone cannot reach
\SI{25}{\percent} — it would need \num{1771} of \num{1615} nodes — and Cumulus in its entirety
holds \SI{3.67}{\percent}.}
\end{table}

The largest single identity holds \SI{10.71}{\percent} of voting weight across 237 Stratus nodes
under both identity models. Decomposed by node key, the three largest hold \num{9480} units across
those 237 nodes (\SI{10.71}{\percent}), \num{9360} across 234 nodes (\SI{10.57}{\percent}), and
\num{3360} across 84 nodes (\SI{3.80}{\percent}), summing to \SI{25.08}{\percent} — that is
\num{22200000}\flux\ of collateral held by three keys. Under the owner-identity model the same
threshold takes four addresses. The Gini coefficient of voting weight is \num{0.8375} across payment
addresses and \num{0.839} across node keys. 47 keys pay to more than one address and 5 addresses
receive from more than one key; no collateral funding transaction pays more than one address
(0 of \num{6178}), so funding-transaction grouping fragments rather than merges and is not by itself
a party model. The fleet is bimodal: 107 keys with ten or more nodes hold \SI{72.5}{\percent} of all
weight, while 334 keys hold exactly one node and are electorally irrelevant. That 819 node keys
serve \num{6489} nodes is a separate observation with two consequences: operators routinely reuse
one node key across many machines, which is a linkage signal here and a key-hygiene matter
elsewhere.

\begin{corollary}[The measured verdict]
\label{cor:23:verdict}
At $\vthresh = 0.25$ and the measured distribution, $\kappa(0.25) = 4$ by payment address --- the
owner-linked proxy, and therefore the operative figure --- and $3$ by node key. Merging addresses
linked by a shared node key or a shared collateral funding transaction into single clusters gives
797 identities and $\kappa(0.25) = 3$ as well. Safety against wrongful eviction is therefore not a
property of the mechanism but the assumption that three or four specific parties do not agree, and
the figure is stable across every identity model computable from public data.
\end{corollary}

Both counts are \textbf{lower bounds} on concentration. They assume every distinct key is a
distinct party; one party holding several keys makes the true number smaller, never larger. They
are not upper bounds on anything. This limitation is structural rather than a gap in the
measurement: a figure that reflects \emph{parties} rather than keys cannot be obtained from public
data.

\begin{corollary}[The graduated tiers are unilateral]
\label{cor:23:tiers}
The proposed graduated-penalty tier at approximately \SI{10}{\percent} is reachable by the largest
single identity alone, which holds \SI{10.71}{\percent} across 237 Stratus nodes under one key. The
tier at approximately \SI{15}{\percent} is reachable by two identities. Any action attached to those
tiers is therefore an action that one or two parties can take at will.
\end{corollary}

\settlednote{Removed --- one threshold, at \SI{25}{\percent} (\Cref{tab:23-recommended}). Actions at the graduated \SI{10}{\percent} and \SI{15}{\percent} tiers, or their removal.
Domain: $\{$none, flag, warn, restrict egress, \ldots$\}$. Trade-off: by \Cref{cor:23:tiers} any
\emph{destructive} action at those tiers is one party's decision; non-destructive, reversible
actions remain consistent with the fail-safe direction. The design intake proposes the tiers without
specifying their effects.}

\begin{corollary}[Raising the threshold does not rescue the property]
\label{cor:23:threshold}
$\kappa \ge 10$ requires $\vthresh \ge 0.40$; $\kappa \ge 16$ requires $\vthresh \ge 0.50$. Even
$\vthresh = 0.51$, the value the design rejected as unreachable in practice, corresponds to
seventeen identities.
\end{corollary}

Raising $\vthresh$ trades liveness (\Cref{prop:23:liveness}) for a larger but still small coalition.
It does not produce a $\kappa$ that could honestly be called a property rather than a promise.

\begin{conjecture}[Custodial dispersion]
\label{conj:23:custody}
If a substantial fraction of the largest owner identities are custodial hosting providers holding
collateral for many distinct beneficial owners, then the number of \emph{parties} behind the four
identities that reach \SI{25}{\percent} is larger than four, and $\kappa(0.25) = 4$ overstates
concentration by however much beneficial ownership is dispersed behind those addresses. Note what
the measurement already rules out: the gain cannot come from separating the node key from the
owner, because the measurement is already taken at the owner --- \texttt{payment\_address} is
derived from the collateral public key
\srccode{fluxd/\allowbreak{}src/\allowbreak{}rpc/\allowbreak{}fluxnode.cpp}{1512--1520} --- and moving between the two identity models
changes $\kappa(0.25)$ only between 3 and 4.
\end{conjecture}

This is a conjecture and not a result, and unlike the other open items in this section it is not
convertible into one: beneficial ownership behind a transparent address is not exposed by any public
data. $\kappa(0.25) = 4$ is therefore the number that stands and the number the paper prints. Note
also that the conjecture cuts both ways — a custodial identity votes with its customers' stake,
which is a different and worse problem than concentration, and it is invisible in the same data.

\notestablished{Open measurement, now narrowed rather than open: the earlier draft of this section
called for concentration by collateral-spending script. That measurement is subsumed, because
\texttt{payment\_address} is derived from the collateral key, so the owner-identity column of the
table above \emph{is} the collateral-key measurement. What remains is beneficial ownership behind
those addresses, which is not computable from public data and which decides
\Cref{conj:23:custody}. No claim about dispersion may be made from public data.}

\settlednote{The start-transaction shape, signed by the collateral (owner) key (\Cref{tab:23-recommended}). Which existing transaction shape the vote follows, and hence which key signs it. Domain:
$\{$the start-transaction shape, signed by the collateral (owner) key
\srccode{fluxd/\allowbreak{}src/\allowbreak{}fluxnode/\allowbreak{}activefluxnode.cpp}{253--288}; the confirm-transaction shape, signed by
the node key \srccode{fluxd/\allowbreak{}src/\allowbreak{}fluxnode/\allowbreak{}activefluxnode.cpp}{290--332}; either, via an on-chain
delegation$\}$. Trade-off: the settled rule says node \emph{owner} keys, which selects the
start-transaction shape --- stake-bound, cold, and the identity every figure in this section is
computed against --- at the cost of asking operators to sign with a cold key on a \SI{24}{\hour}
deadline. The confirm-transaction shape is operationally far easier because that key is already hot
and already signs every 500 blocks, but it is the VPS key, it may be custodial, and choosing it
would require recomputing the entire concentration measurement against the node-key column and
restating $\kappa(0.25)$ as 3.}

% =====================================================================
\subsection{The negative result}
\label{sec:23:negative}

\begin{remark}[The measured concentration is accepted, not disputed]
\label{rem:23:accepted}
Before the negative result, the team's position on it. Presented with the measurement --- four owner
identities reaching the \SI{25}{\percent} threshold, three under the strongest linkage computable
from public data --- the response was to accept it: if four parties hold a quarter of the network's
collateral then those parties should be able to remove an application, and ``as network grows it will
be more decentralized'' \srcintent{08-answers-round-2.md, round 11} (\planned).

The paper records this as a deliberate acceptance rather than an unrecognised risk, which changes how
the rest of this section should be read: what follows is not an objection the design has failed to
answer, but a property the design has been told about and kept. Two observations remain load-bearing
even so, and neither is a criticism of the choice.

\textbf{First, growth dilutes concentration only if new collateral is new parties.} The measurement
is a ratio, not a count. If the network doubles and the four largest owners double with it, the ratio
is unchanged; dilution requires that marginal entrants be new and independent, which is an empirical
question this corpus cannot settle either way. \textbf{Second, and more sharply, the mechanism the
paper elsewhere identifies as the response to weak demand runs the other way.}
\Cref{rem:pnr:contraction} records that a demand shortfall is answered by contraction --- operators
exit, the survivors earn more --- and contraction concentrates. Growth as a remedy and contraction as
a stabiliser cannot both be relied on; which one obtains depends on demand for Flux Cloud, and no part
of this section's design references demand at all.
\end{remark}

The natural response to \Cref{cor:23:verdict} is to change the weight function. It does not work,
and the reason it does not work is worth stating precisely, because it is the section's
contribution.

Capping per-identity weight at some $c$ is defeated at negligible cost: an address is free, each
node's collateral can be respent to its own, and the 237-node identity becomes 237 identities. The
cost of that split is one collateral respend, 100 blocks of maturity and one start transaction per
node \srccode{fluxd/\allowbreak{}src/\allowbreak{}coins.cpp}{630--686} — real, but bounded by a few hours of missed rotation
against the \num{22128625}\flux\ the threshold requires, so it is not a barrier to anyone who can
reach the threshold in the first place. This is exactly the Sybil that
collateral weighting was chosen to prevent, reintroduced with the cap itself as the incentive.
Quadratic weighting per identity is defeated identically and more sharply — the same holder is
worth $\sqrt{9480} \approx 97.4$ as one identity and $237\sqrt{40} \approx 1498.9$ as 237, a
fifteenfold reward for splitting. Only weightings applied \emph{per node} are Sybil-neutral, and
those are tier reweightings, not concentration fixes:

\begin{table}[H]
\centering
\small
\begin{tabular}{@{}lrrrrrrr@{}}
\toprule
weighting & Stratus & top key & top 3 & $\kappa(0.25)$ & $\kappa(0.51)$ & Gini
  & cost of \SI{25}{\percent} (\flux) \\
\midrule
collateral (settled) & \SI{73.5}{\percent} & \SI{10.71}{\percent} & \SI{25.08}{\percent}
  & \textbf{3} & 17 & \num{0.839} & \num{22128625} \\
tier base $1:3.5:9$ & \SI{62.2}{\percent} & \SI{9.06}{\percent} & \SI{21.23}{\percent}
  & 5 & 22 & \num{0.793} & --- \\
$\sqrt{\coll{t}/\num{1000}}$ per node & \SI{53.5}{\percent} & \SI{7.79}{\percent}
  & \SI{18.31}{\percent} & 6 & 26 & \num{0.770} & --- \\
one node, one vote & \SI{25.1}{\percent} & \SI{7.40}{\percent} & \SI{17.58}{\percent}
  & 6 & 38 & \num{0.716} & \num{1623000} \\
\bottomrule
\end{tabular}
\caption{Sybil-neutral reweightings, computed over the same \num{6489} records
\srcmeasured{api.runonflux.io/\allowbreak{}daemon/\allowbreak{}listzelnodes}{2026-09-03}. Every row costs
\num{1000}\flux\ to add one Cumulus node, so none reopens the Sybil vector; the last column is the
capital needed to \emph{buy} the threshold outright, and only the two extreme rows were computed.
Dashes are uncomputed, not zero.}
\end{table}

\begin{remark}[Which identity model this table uses, and how to read it]
\label{rem:23-table-identity}
Every row above is computed under the \textbf{node-key} model, in which the settled weighting gives
$\kappa(0.25) = 3$ rather than the operative $4$ that the owner-identity model yields. The
owner-identity recomputation has not been run for the per-row figures.

This is a caveat about reading individual rows, not a defect in the conclusion. The recomputation
would shift each row by the same one-or-two identities that \Cref{sec:23:measurement} measures between
the two models --- it moves every row in the same direction by roughly the same amount, which is why
the open problem below is stated as \emph{no row approaches a $\kappa$ that could be called a
property} rather than as a claim about any particular row. The aggregate conclusion is robust to the
model choice; a single row quoted on its own is not, and should be recomputed under the owner-identity
model first.
\end{remark}

\begin{openproblem}[No reweighting, threshold or listed remedy raises $\kappa$ meaningfully]
\label{op:23:negative}
At the measured distribution and $\vthresh = 0.25$: no weight function in the Sybil-neutral design
space raises $\kappa(0.25)$ above 6, and the best of them --- one node, one vote --- reduces the
capital required to buy the threshold from \num{22128625}\flux\ to \num{1623000}\flux, a factor of
\num{13.6}, which is a strictly worse trade. Raising $\vthresh$ to \num{0.51} yields
$\kappa = 17$ at a liveness cost the design has already rejected (\Cref{cor:23:threshold}). None of
the remedies listed in \Cref{sec:23:remedies} changes $\kappa$ at all, and neither does the on-chain
shape of \Cref{sec:23:mechanism}: moving the tally into consensus changes who can \emph{verify} the
decision, not who can \emph{make} it. \emph{Open:} whether any
mechanism consistent with the settled decisions --- collateral weighting, absolute threshold,
abstention-as-no, no collateral slashing --- achieves a $\kappa^{\star}$ at $\vthresh = 0.25$
materially above the measured 4 (3 under the node-key model)
without reintroducing a Sybil vector. This paper does not have one, and the two-chamber rule of
\Cref{sec:23:remedies} explicitly does not claim to be one.
\end{openproblem}

The reason is not subtle once stated: concentration lives in node \emph{ownership}, not in the
weight function. The largest key holds \SI{7.4}{\percent} of nodes by count, so even one-node-one-vote
leaves it above the \SI{5}{\percent} mark; the weight function only decides how much worse than
that the ownership distribution is allowed to make things.

% =====================================================================
\subsection{Dilution: why the honesty incentive does not bind}
\label{sec:23:dilution}

The design's argument for why free voting is not a vulnerability is that operators are paid from
application revenue, so evicting applications reduces their own income and they are therefore
incentivised to be honest \srcintent{evidence/design/04-moderation.md}. The argument is correct in
aggregate and wrong at the margin, which is the standard shape of a collective-action problem. It
fails in three steps, each independently sufficient.

\subsubsection{The current term is already paid}

\begin{proposition}[Eviction is income-neutral to the fleet for the current term]
\label{prop:23:neutral}
Under \PNR{} as specified in \Cref{sec:07-pnr}, hosting is prepaid for $\dur$ blocks and the price
$\price$ is split by consensus at payment time: $\foundshare\price$ to the Foundation and
$\nodeshare\price$ into the pool $\pool$, which pays all nodes a drip of $\pool/\drip$ per block
regardless of which applications run. Eviction of $\app$ leaves $\pool$ unchanged. Therefore the
fleet's income for the current term is invariant under eviction, and for a node currently hosting
$\app$ the effect is strictly positive: unchanged pay, released resources.
\end{proposition}

\begin{proof}
$\pool$ is credited at payment and drained by a per-block drip that is a function of $\pool$ and
$\drip$ only; no term of the drip references the set of running applications. Removing $\app$ alters
no argument of that function. A hosting node's payment is likewise independent of hosting, so its
income is unchanged while its committed CPU, RAM and disk are freed.
\end{proof}

What the fleet loses is the \emph{future}: renewals of $\app$, and the demand effect of a network
where eviction is capricious. Let $R_\app$ be expected future renewal revenue in \flux, $g_i \ge 0$
the resource relief to a hosting node, and $b_i \ge 0$ any private benefit to identity $i$ from
$\app$'s removal. With $s_i$ node $i$'s share of the pool, node $i$ prefers eviction iff
$b_i + g_i > \nodeshare\, s_i\, R_\app$.

\subsubsection{Dilution}

\PNR{} splits the drip across tiers in the ratio $1 : 3.5 : 9$ and then equally within a tier, so
$s_i = (\tbase{t_i}/13.5)/\ncount{t_i}$:

\begin{table}[H]
\centering
\small
\begin{tabular}{@{}lrrr@{}}
\toprule
tier & $s_i$ (\PNR{} tier-base split) & $s_i$ (collateral share) & cost of evicting $\app$ at $\nodeshare = 0.20$ \\
\midrule
Cumulus & \SI{0.0023}{\percent} & \SI{0.0011}{\percent} & \SI{0.00046}{\percent} of $R_\app$ \\
Nimbus  & \SI{0.0161}{\percent} & \SI{0.0141}{\percent} & \SI{0.0032}{\percent} of $R_\app$ \\
Stratus & \textbf{\SI{0.0410}{\percent}} & \textbf{\SI{0.0452}{\percent}} & \textbf{\SI{0.0082}{\percent}} of $R_\app$ \\
\bottomrule
\end{tabular}
\caption{Per-node pool share and the private cost of destroying an application's future revenue.}
\end{table}

A bloc of 554 Stratus nodes — the \SI{25}{\percent} threshold — internalises
$554/\num{1627} \times 9/13.5 = \SI{22.7}{\percent}$ of the operator share, i.e.\
\SI{4.54}{\percent} of $R_\app$ at $\nodeshare = 0.20$. The largest single identity internalises
\SI{1.94}{\percent}. The social cost of a wrong eviction is $\nodeshare R_\app$ to operators, plus
$\foundshare R_\app$ to the Foundation, plus the tenant's unused prepayment, plus the loss to the
application's users; the decisive coalition bears under \SI{5}{\percent} of the first term alone.

This is a collective-action problem in textbook form: the fleet collectively wants applications and
no individual member is deterred from destroying one. The two design decisions that produce it are
each correct on their own terms. Paying \emph{all} nodes is the fair rule — placement is
luck-driven and over a rotation every node is paid equally — and it is precisely what makes
$s_i \approx 4 \times 10^{-4}$. Paying only hosting nodes would concentrate the deterrent where the
workload sits, and would reintroduce exactly the placement-metering dependency \PNR{} was designed
to avoid. The team should see the interaction plainly: \textbf{fairness in \PNR{} and honesty in
moderation pull on the same rope in opposite directions.}

\subsubsection{The deterrent is weakest at launch}

\begin{proposition}[Launch-time deterrent]
\label{prop:23:launch}
At \PNR{} activation, with $\nodeshare = 0.20$, operator emission of \num{6386040}\flux\ per year
and application revenue at approximately \SI{4.3}{\percent} of operator income, the expected future
revenue a Stratus operator forgoes by destroying one average application is approximately
\num{0.10}\flux\ per year.
\end{proposition}

\begin{proof}[Proof sketch]
Application revenue at \SI{4.3}{\percent} of total operator income against \num{6386040}\flux\ of
emission gives \num{286320}\flux\ per year from all applications (\Cref{sec:07-pnr}). The Stratus
tier receives $9/13.5$ of the operator share, i.e.\ about \num{190880}\flux\ per year, spread over
\num{1627} nodes, i.e.\ about \num{117.3}\flux\ per node per year. Divided across the \num{1193} priced
applications in the \PNR{} demand baseline \srcintent{plan/mech-pnr.md} — the registry returns
\num{1195} specifications
\srcmeasured{api.runonflux.io/\allowbreak{}apps/\allowbreak{}globalappsspecifications}{2026-09-03} — that is
\num{0.10}\flux\ per Stratus node per application per year.
\end{proof}

Against a forfeiture floor of \num{0.01}\flux\ on the tenant's side, that is the entire economic
content of the honesty argument at launch. It rises with the ramp toward $\nodesharecap = 0.50$ and
with demand, and it is smallest on the day the mechanism would ship. Note also the self-dealing
case, which is worse than neutral: an operator who also hosts $\app$ has $g_i > 0$ by
\Cref{prop:23:neutral}, so under \PNR{} there is a small positive incentive to evict one's own
tenants, bounded only by the threshold.

Finally, a constraint this section exports to the revenue mechanism. \Cref{sec:07-pnr} must
guarantee that \textbf{forfeitures from eviction never enter the fee pool}. If they did, evicting an
application would pay the fleet that voted to evict it, converting a socialised cost into a
distributed gain and inverting the incentive this subsection has just shown to be weak.

\settlednote{Burn pro-rata (\Cref{tab:23-recommended}, \Cref{rem:o16-burn}). Disposition of the pool's unearned remainder on eviction. Domain: $\{$keep in $\pool$
(status quo); burn pro-rata$\}$. Trade-off: burning the unearned fraction is the only construction
identified that gives voters a non-zero \emph{current-term} cost of eviction and thereby repairs
\Cref{prop:23:neutral}; it requires \PNR{} payments to be bound to individual applications, which
the private-payment section notes may leak the workload under shielded payment. Keeping the
remainder preserves \Cref{prop:23:neutral} exactly as measured.}

% =====================================================================
\subsection{Remedies, evaluated}
\label{sec:23:remedies}

Each remedy below is assessed on three axes: what it fixes, what it costs, and what it does not
fix. The first is adopted by the specified design and the rest are not; the third column is the
same in every one of them, and it is the actual problem.

\paragraph{Proposal capacity bounded by collateral (adopted, \Cref{def:23:capacity}).}
\emph{Fixes:} report spam, at a price of \num{1000}\flux\ of locked collateral per concurrently open
report; block-space exhaustion by fee-less transactions, which nothing else in the design bounds;
and, as a side effect, reporter eligibility, which stops being a parameter. \emph{Costs:} nothing
new — no bond, no custody script, no adjudicator that must learn the outcome to release funds, and
no forfeiture destination to argue about. It is a validity predicate over data \fluxd{} already
holds. \emph{Does not fix:} capacity is distributed as node counts, so the ability to occupy the
report channel is concentrated in the same hands as everything else (Gini \num{0.716}); and it
prices nuisance in \emph{locked} rather than \emph{spent} capital, so a determined griefer pays
only opportunity cost. Concentration.

\paragraph{Refundable report bond (evaluated, not adopted).} \emph{Would fix:} frivolous reports,
by making $n$ failed reports cost $n\bond$ outright rather than $n$ node-slots of locked capital.
\emph{Costs:} a refundable bond needs an adjudicator that can read the outcome — a consensus change,
a release key, or a non-refundable fee that taxes honest reporters — and each of those is a new
trusted or new consensus object. The capacity rule obtains a comparable bound with none of them,
which is why the design has no bond and why $\bond$ appears nowhere in this section's arithmetic.
\emph{Does not fix:} any report that succeeds is free either way, so a coalition at threshold pays
nothing. Concentration.

\paragraph{Rate limiting.} \emph{Fixes:} denial of service against operator attention; serial
re-reporting of one tenant. \emph{Costs:} per-key limits multiply with keys — a 480-node operator
who gives each node its own key obtains $480\,r_{\max}$ at zero collateral cost, and per-node limits
are the same thing. This is precisely the failure the capacity rule avoids by counting
\emph{nodes} rather than \emph{keys}: by \Cref{prop:23:capacity} the total is invariant under any
re-keying, so the rule is a genuine bound where a per-key rate limit is not. \emph{Does not fix:}
anything about who can reach threshold. Concentration.

\paragraph{Eligibility gating.} \emph{Fixes:} excludes dead and unbenchmarked nodes, and — because
confirmation in $N(\hgt)$ already requires a benchmark-bearing confirmation transaction within the
last 640 blocks — makes ``confirmed at $\hgt_0$'' the on-chain proxy for benchmark-passing, which is
also what keeps the tally verifiable against consensus state rather than against a
\FluxOS{}-local status. \emph{Costs:} the benchmark signature is not operator-bound, as the
system-model section establishes, so the gate is only as strong as the benchmark; any stricter
``good standing'' test that is not on-chain breaks verifiability. \emph{Does not fix:} the largest
keys are benchmark-passing nodes in good standing. Concentration.

\paragraph{Reinstatement (rejected: the team's decision is that a ban is final;
\Cref{op:23:reinstatement}).} \emph{Would fix:}
a symmetric reinstatement quorum gives an honest majority a path to reverse a wrongful ban; a
\emph{sunset}, under which a ban expires unless renewed by a fresh report and a fresh quorum, puts
the liveness burden on the side that has already demonstrated it can reach quorum — a persistent bad
application is re-banned each period, a wrongly banned one returns by default. \emph{Costs:} a
symmetric quorum inherits \Cref{lem:23:turnout} and is asymmetric in
exactly the wrong direction, since the dispersed honest majority faces the \SIrange{33}{50}{\percent}
turnout problem while the coalition that banned can reinstate at will; a sunset gives a harmful
application a periodic window of uptime, costs the banning side a renewed quorum every period, and
converts the ban set from monotone (\Cref{prop:23:monotone}) into state that must be re-derived at
every height. \emph{Does not fix:} under a symmetric quorum the same four identities decide both
directions. The specified design adopts neither, consistent with the team's decision that a ban is
final and the remedy is redeployment (\Cref{op:23:reinstatement}).
Concentration.

\paragraph{Opt-in delegation.} \emph{Fixes:} liveness. Turnout rises without operators following
every report, and a higher, safer $\vthresh$ becomes reachable later. \emph{Costs:} delegation
concentrates weight in delegates by construction, which is the documented trajectory of every
delegated token-vote system, and makes delegates the target of the coalition. \emph{Does not fix:}
concentration — it measurably worsens it.

\settlednote{Evidence is an optional content hash, and a report counts as its reporter's vote at $\hgt_0$ (\Cref{tab:23-recommended}). The report's descriptive fields, and whether a report is itself a vote. Domains:
$\mathsf{evidence} \in \{$absent; optional; a mandatory content hash$\}$; $\mathsf{category}$ as
left open in \Cref{sec:23:c5}; and whether $\mathrm{own}_r$'s own weight is added to $Y_r$ at
$\hgt_0$ or requires a separate vote transaction. Trade-offs: an evidence hash is unverifiable by
any node and costs bytes in a fee-less transaction, and its only real consumer is a future
adjudication of the kind the team's fallback contemplates; counting the report as a vote saves one
transaction and makes the report's author's weight implicit rather than explicit in the chain
record, which is a small loss of the attribution the design otherwise provides. Note that three
parameters the earlier draft carried here are now \emph{settled} by the on-chain shape rather than
chosen: the per-application cooldown is zero, the network-wide cap on open reports is
$\lvert N(\hgt) \rvert$ (\Cref{cor:23:capacity}), and reporter eligibility is $I(\hgt_0)$.}

\subsubsection{The recommended construction, and exactly what it buys}

\begin{definition}[Two-chamber decision rule]
\label{def:23:twochamber}
Let $\mathrm{tenure}_i(\hgt) \defeq \hgt - \min_{j\,:\,\mathrm{own}_j = i}
\mathrm{added\_height}_j$, the age of the identity's oldest node. Carry the report iff \emph{both}
\[
Y_r \;\ge\; \vthresh\,\vtotal^{\!\star}(\hgt_0)
\qquad\text{and}\qquad
\bigl\lvert\{\, i \in V_r \ :\ \mathrm{tenure}_i(\hgt_0) \ge T \,\}\bigr\rvert \;\ge\; m,
\]
with $m \in \mathbb{N}$ and $T \in \mathbb{N}$ (blocks) open.
\end{definition}

\begin{proposition}[The two-chamber rule preserves everything settled and lowers nothing]
\label{prop:23:twochamber}
Adding the second conjunct leaves weights, threshold, abstention rule and enforcement untouched;
\Cref{lem:23:turnout}, \Cref{thm:23:turnout} and \Cref{thm:23:collusion} continue to hold verbatim
for the first chamber; the vote transactions are unchanged, with validation extended by a
tenure lookup against $N(\hgt_0)$, which is data \fluxd{} already holds. In particular
$\kappa(0.25) = 4$ is unchanged.
\end{proposition}

\begin{proof}
The rule is a conjunction, so the set of evictable configurations shrinks weakly; every statement of
the form ``no set smaller than $X$ can evict'' is preserved and no statement of that form is
strengthened, because the second conjunct is satisfiable by any coalition that has positioned $m$
tenured identities. \Cref{lem:23:turnout} and \Cref{thm:23:turnout} concern the first conjunct only
and are untouched.
\end{proof}

The tenure gate is not optional. Without it the second chamber counts identities, and an identity is
nearly free: a four-identity coalition re-addresses $m-4$ of its own nodes and satisfies the chamber
at no collateral cost at all. Re-addressing requires a collateral respend and a new start
transaction and, \srcinferred, resets \texttt{added\_height} and the node's position in the payment
rotation — one missed Stratus rotation
is \num{9}\flux, about \SI{0.155}{\percent} of that node's annual subsidy income. That inference
is load-bearing for the parameter choice and is not verified in this survey; it must be confirmed in
\fluxd{} before $m$ and $T$ are set.

Tenure eligibility at the measured distribution:

\begin{table}[H]
\centering
\small
\begin{tabular}{@{}lrrrr@{}}
\toprule
$T$ & nodes with tenure $\ge T$ & distinct keys & tenured weight & top-3 keys' tenured weight \\
\midrule
\SI{7}{\day} (\num{20160} blocks)   & \num{5639} (\SI{86.9}{\percent}) & 744 & \SI{90.1}{\percent} & \SI{24.05}{\percent} \\
\SI{30}{\day} (\num{86400})         & \num{4149} (\SI{63.9}{\percent}) & 607 & \SI{77.0}{\percent} & \SI{23.95}{\percent} \\
\SI{90}{\day} (\num{259200})        & \num{2824} (\SI{43.5}{\percent}) & 412 & \SI{57.2}{\percent} & \SI{19.66}{\percent} \\
\SI{180}{\day} (\num{518400})       & \num{1516} (\SI{23.4}{\percent}) & 286 & \SI{31.5}{\percent} & \SI{9.49}{\percent} \\
\bottomrule
\end{tabular}
\caption{Tenure distribution \srcmeasured{api.runonflux.io/\allowbreak{}daemon/\allowbreak{}listzelnodes}{2026-09-03}. The
last column is the weight, as a share of all weight, of the three largest keys among those with
tenure $\ge T$; it is informational — the tenure gate applies to the identity chamber only — but
that the three largest keys' nodes are mostly younger than \SI{180}{\day} says the concentration is
recent.}
\end{table}

What the rule buys is precise and limited. A reactive eviction — ``these four identities sign now'' —
would no longer be possible unless $m-4$ other tenured identities also sign. A \emph{pre-planned}
eviction would remain possible for any coalition that positioned $m - \kappa(\vthresh)$ tenured
identities at least $T$ blocks earlier, at a cost of $(m-4) \times \num{1000}\flux$ in Cumulus
collateral, or nothing at all by re-addressing nodes it already owns, and in either case
\textbf{visible on-chain in advance}. At the settled
$m = 21$ the collateral route costs \num{17000}\flux, which is \SI{0.077}{\percent} of the
\num{22100000}\flux\ the first chamber already requires; at $m = 100$, the top of the domain, it
costs \num{96000}\flux, \SI{0.43}{\percent}. The rule is therefore not a capital
barrier. It is a time-and-visibility barrier, and its entire value lies in composition with
publishing $\kappa(\vthresh)$, the top-$N$ shares and the Gini at every reduction height: a burst of
re-addressing, or of new single-node identities from a large operator, becomes a public, datable
warning.

It costs liveness in the second chamber — $m$ distinct tenured identities must actually vote. With
the settled $T = \SI{30}{\day}$ and $m = 21$, \SI{3.5}{\percent} of eligible keys must participate
(\SI{16.5}{\percent} at $m = 100$); with $T = \SI{90}{\day}$, \SI{5.1}{\percent}. Median node
tenure is \num{193620} blocks, about
\SI{67}{\day}, so any $T$ beyond roughly \SI{60}{\day} disenfranchises the median node from the
second chamber.

\settlednote{$m = \num{21}$, $T = \num{86400}$ blocks (\Cref{tab:23-recommended}). $m$, the identity-chamber count, and $T$, the identity tenure. Domains:
$m \in [5,\ 100]$; $T \in [\num{20160}\ (\SI{7}{\day}),\ \num{259200}\ (\SI{90}{\day})]$ blocks.
Trade-offs: larger $m$ raises the Sybil cost of pre-positioning to $(m-\kappa)\times\num{1000}\flux$
or to $T$ blocks of waiting, and raises the number of tenured identities that must actually vote;
larger $T$ raises the attacker's required lead time and disenfranchises younger honest identities
against a median tenure of \SI{67}{\day}. The second chamber counts the same owner identities as the
first (\Cref{def:23:electorate}); it is the tenure gate, not the identity definition, that does the
work, because every identity definition available here is cheap to multiply.}

This is the honest ceiling. Raising the cost of \emph{surprise} rather than the cost of
\emph{capture} is what mechanism design can do here, and the section says so rather than presenting
a bicameral rule as a safety proof it is not.

% =====================================================================
\subsection{Attacks and failure surface}
\label{sec:23:attacks}

\paragraph{Report brigading.} $n$ colluding reporters filing $n$ reports on one application would
achieve nothing beyond one report: at most one report per application hash is open, and the tally is
per identity, not per report. Brigading the \emph{vote} is simply the coalition of
\Cref{thm:23:collusion}; it costs nothing at all, because there is no bond and no fee. Brigading the
\emph{report channel} is bounded by \Cref{prop:23:capacity} at one concurrent report per node
operated, which is the whole of the anti-spam design. Cost to the network per brigade: one fee-less
transaction per report, plus whatever vote traffic the reports attract, against the block-space
budget computed in \Cref{sec:23:mechanism}.

\paragraph{A competitor evicting a rival.} Requires $\vthresh\vtotal^{\!\star}$ of weight. If the
competitor is one of the four largest owner identities, the other three suffice with it. If not, it
must acquire \num{22100000}\flux\ of collateral, wait 100 blocks of maturity plus confirmation,
and — under the two-chamber rule — pre-position tenure; or it must rent the signatures. Its payoff
is the rival's diverted demand less $\nodeshare s_i R_\app$, which for the largest identity is about
\SI{2}{\percent} of $R_\app$. Nothing is staked and nothing is forfeited on either outcome. The
record is public, so the deterrent is reputational and, under the fallback below, retrospective.

\paragraph{A censoring bloc.} Four owner identities — three under the strongest linkage computable
from public data — would ban any application at will within
$\Delta_{\mathrm{vote}} + \delta + \Delta_{\mathrm{exec}}$ blocks at zero cost, and at
\SI{50}{\percent} (nineteen identities) they would satisfy any two-chamber rule they had prepared
for. Nothing in the mechanism prevents this, and under \Cref{op:23:reinstatement} nothing reverses
it: a wrongful ban is a permanent chain fact. What the mechanism does is make it
\emph{attributable}, and the on-chain shape strengthens exactly this: the evidence is not a
gossiped certificate that a node may or may not hold but a set of confirmed, height-stamped
transactions that every archival node has, in an order the chain fixes.

\paragraph{Producer-level censorship.} New with the on-chain shape, and the reason the inclusion
rule of \Cref{prop:23:inclusion-required} was made consensus-required rather than left to producer
discretion. Under the settled rule a producer cannot decline to include a due ban memo and still
produce a valid block, so this vector is closed by a consensus rule rather than merely made
statistically unlikely: the delay is one slot for a report with no carried report queued ahead of it,
regardless of $f$, the fraction of nodes that would have preferred to withhold it
(\Cref{prop:23:inclusion}). What remains subject to producer
discretion is \emph{vote} censorship: a coalition holding a fraction $f$ of nodes can decline to
include a given vote, but it needs every one of \num{2880} slots to succeed, so vote censorship
fails with probability $1 - f^{\num{2880}}$. Putting the decision on-chain now makes both the vote
and the memo robust to an individual producer --- the vote by the size of the window, the memo by an
explicit consensus rule --- which is the reverse of the asymmetry an earlier draft of this section
anticipated.

\paragraph{Electorate timing.} The reporter chooses $\hgt_0$. Without the trailing maximum of
\Cref{def:23:electorate}, a reporter would choose a height at which many nodes are temporarily
unconfirmed. The largest two keys hold \SI{21.28}{\percent} and reach \SI{25}{\percent} of any
$\vtotal(\hgt_0) \le 0.851\,\vtotal$; the largest alone reaches it if
$\vtotal(\hgt_0) \le 0.428\,\vtotal$. Confirmation expiry is 640 blocks, so an incident that leaves
\SI{15}{\percent} of weight unconfirmed for a few hours is the attack window, and $\Delta_W$ closes
it. Racing a vote to the window boundary is not an attack — a vote confirming after
$\hgt_0 + \Delta_{\mathrm{vote}}$ is simply invalid — but a reorg across the ban block is, and
$\Delta_{\mathrm{exec}}$ absorbs it.

\paragraph{Griefing, and what actually bounds it.} The specified design's answer is
\Cref{def:23:capacity} and nothing else: a griefer's sustained report rate is its node count per
\SI{24}{\hour}, and its concurrent footprint is its node count, priced at \num{1000}\flux\ of
locked and recoverable collateral per slot. Holding every deployed application under permanent
report therefore costs \num{1195000}\flux\ of locked capital and buys operator attention rather than
eviction. Compare the alternatives the design did not take: a bond would charge per \emph{failed}
report and nothing to a griefer with a whale's backing; a per-key rate limit would be multiplied by
free keys, since a 480-node operator has 480 of them; eligibility gating alone requires
\num{1000}\flux, which is already true of every operator. Under the two-chamber rule, a griefer
who reaches the first chamber must also find $m$ tenured identities.

\paragraph{Wash-deployment.} A deployer who registers and then reports their own application spends
$\price$ and, on success, forfeits the unused remainder of it; on expiry, spends nothing further,
because reporting is free and the only cost was the capacity slot. There is no state in which the
cycle yields more than it costs, because the mechanism has no pool a wash could farm — which is
exactly the property that the \Cref{sec:07-pnr} constraint on forfeitures preserves.

\paragraph{The quorum is simply wrong.} An honest coalition bans a legitimate application on
mistaken evidence. Recourse, in order of availability: redeployment under a new name, always
available, costing one fee and preserving none of the old instance's reachability, since \FDM{} and
\PDNS{} names are keyed per application name; and the public chain record, which names the voters.
Reinstatement is \emph{not} on that list, because the design has none
(\Cref{op:23:reinstatement}). When redeployment does not suffice, \textbf{the design bounds the
damage rather than preventing the error}: the maximum loss is one prepaid term (floor
\num{0.01}\flux, typically the unused fraction of $\price$), the redeployment fee, and
$\Delta_{\mathrm{vote}} + \delta + \Delta_{\mathrm{exec}}$ blocks plus redeployment time of
downtime. That bound is the ``err toward leaving content up'' dial applied to the cost of the other
error, and it is the design's actual guarantee. It is also the whole of the guarantee: the ban
itself is permanent.

\paragraph{Byzantine minority and colluding majority.} A minority below $\vthresh$ can report
(bounded by its node count), can vote, cannot ban, cannot forge or replay votes, cannot keep a vote
out of the chain for a whole window (\Cref{prop:23:inclusion}), and cannot produce a valid ban memo,
because a block carrying one for an uncrossed report is invalid and is rejected like any other
invalid block. A coalition at or above $\vthresh$ is the censoring bloc above. Sustained low demand
shrinks $\pool$ and drives the already-negligible \PNR{} deterrent toward zero without affecting
safety. Sustained low participation means genuinely harmful applications stay up, which is a
liveness failure and not a safety failure (\Cref{prop:23:expiry}).

\paragraph{The team's stated fallback.} If malicious reporting appears, the chain is to be forked to
introduce penalties for malicious submissions \srcintent{evidence/design/04-moderation.md}. This is
a legitimate governance answer and the paper states it plainly, with its limits. It is
\emph{reactive}, acting after the ban. It requires a \emph{hard fork} of \fluxd, with the
coordination that implies — and the consensus section's history of what a consensus change costs
this network is directly relevant. It requires \emph{distinguishing malicious from mistaken reports
after the fact}, for which no oracle exists and which the mechanism does not attempt. And because
collateral is never slashed, the penalty must be something other than collateral, which is not
specified. What the mechanism does contribute is the record: a future fork would have signed,
height-stamped, chain-resident evidence of every vote to act on, and would not have to trust anyone
to have kept it. Without that the fallback would have nothing. It is worth noting that a fork is
also the only route the specified design leaves for undoing a wrongful ban, which makes
\Cref{op:23:reinstatement} and this paragraph the same problem seen from two sides.

% =====================================================================
\subsection{The decentralization endgame}
\label{sec:23:endgame}

A section on governance owes the reader an inventory of what is team-operated today and must not be
at target. Drawing on the networking and addressing section's centralization table and on the
node-services extraction, the components are: \texttt{api.\allowbreak{}runonflux.\allowbreak{}io} as the sole
application-discovery source, with no on-chain fallback and no independent replica
\srccode{flux-domain-manager/\allowbreak{}src/\allowbreak{}services/\allowbreak{}flux/\allowbreak{}dataFetcher.js}{47--91}; the Foundation-operated
private \texttt{10.100.0.0/24} carrying the WireGuard hub, the DNS gateway, the Arcane load
balancers and both PowerDNS masters
\srccode{flux-fdm-infrastructure/\allowbreak{}ansible/\allowbreak{}inventory/\allowbreak{}hosts.yaml}{102--206}; the 16-host \FDM{} fleet
\srccode{flux-fdm-infrastructure/\allowbreak{}ansible/\allowbreak{}inventory/\allowbreak{}hosts.yaml}{8--100},
none of which runs on a node operator's machine; the \texttt{fluxos-\allowbreak{}network-\allowbreak{}policy} channel of
\Cref{def:23:policy-channel}; the single-instance \texttt{fluxstorage} service, which holds
oversized environment variables and commands unencrypted at rest with no write-side protection at
all \srccode{fluxstorage/\allowbreak{}src/\allowbreak{}routes.js}{17--67}; the release API, which performs no cryptographic
verification of what it serves as the latest release
\srccode{fluxos-release-api/\allowbreak{}src/\allowbreak{}flux\_release\_api/\allowbreak{}api.py}{176--229}; the Foundation-held application
owner keys used for auto-renewal \srccode{appsmonitor/\allowbreak{}config/\allowbreak{}default.js}{1--3}; and a
\texttt{fluxteam} privilege level in the node UI strictly above a node owner's own \texttt{admin}
\srccode{fluxos-frontend/\allowbreak{}src/\allowbreak{}navigation/\allowbreak{}vertical/\allowbreak{}fluxadmin.js}{39--43}. To that list the consensus
layer adds the 2-of-4 emergency-block key set \srccode{fluxd/\allowbreak{}src/\allowbreak{}emergencyblock.cpp}{24--30,183--190} —
3-of-4 with fresh keys intended, aligning with the bridge quorum (\S22), unscheduled, \planned
\srcintent{evidence/design/08-answers-round-2.md, round 5 answers} — the
consensus-enforced dev-fund remainder paid to a hardcoded address
\srccode{fluxd/\allowbreak{}src/\allowbreak{}main.cpp}{2877--2884}, and the single transparent address to which all
application payments land today by policy
\srcmeasured{api.runonflux.io/\allowbreak{}apps/\allowbreak{}deploymentinformation}{2026-09-03}.

One further fact belongs here rather than being left for a reader to discover. The largest in-flight
change to how applications are described on Flux — the \specv{9} work — currently lives on an
individual contributor's fork rather than in the organisation's repository: \num{1167} commits ahead
of the organisation's development branch, and 932
changed files with $+\num{180317}/-\num{49561}$ lines against upstream master, actively pushed
\srcmeasured{gh api repos/\allowbreak{}RunOnFlux/\allowbreak{}flux/\allowbreak{}compare/\allowbreak{}development...MorningLightMountain713:feat/\allowbreak{}fluxos-v9}{2026-09-03}.
That is an ordinary open-source pattern and not a criticism. For a paper arguing progressive
decentralization it is nonetheless a fact about where development authority sits, and the paper
states it rather than omitting it.

The order in which these move matters, and one ordering constraint is already visible: the policy
channel's own stated next step — signed documents with a rollback-resistant sequence number
\srccode{fluxos-network-policy/\allowbreak{}README.md}{323--326} — is strictly cheaper than the quorum, requires
no consensus change, and would raise the deployed mechanism from ``no authentication'' to
``authenticated by a named key''. It would not raise $\kappa_{\mathrm{deployed}}$ above 1. A
threshold over publisher keys would, and is the obvious intermediate step between
\Cref{prop:23:kappa-today} and the quorum.

% =====================================================================
\subsection{Recommended values for the open parameters}
\label{sec:23:recommended}

Every parameter this section left open is collected here with a value. Each was proposed by this
paper and \textbf{confirmed by the author} \srcintent{08-answers-round-2.md, round 11}; the reasoning
is retained so that any of them can be overridden with a number rather than re-argued from the
beginning. Block counts assume the \SI{30}{\second}
post-\PoN{} target, i.e.\ \num{2880} blocks per day \srccode{fluxd/src/chainparams.cpp}{105}.

\begin{table}[H]
\centering\footnotesize
\caption{Parameter values for the content-reporting quorum. Proposed by this paper and
\textbf{confirmed by the author} \srcintent{08-answers-round-2.md, round 11}. \planned,
\srcinferred{} for the derivations except where a citation is given.}
\label{tab:23-recommended}
\begin{tabular}{@{}llp{0.44\linewidth}@{}}
\toprule
Parameter & \textbf{Recommended} & Reason \\
\midrule
$\Delta_W$ trailing window & \textbf{\num{20160}} (\SI{7}{\day}) & trailing \emph{maximum}; an
attacker must wait a week after any contraction, and staleness fails safe \\
$\mathsf{app\_hash}$ scope & \textbf{$(\mathsf{app\_name}, \mathsf{owner})$} & a spec-version hash is
evaded by re-broadcasting a trivially edited spec \\
$\mathsf{decision}$ field & \textbf{evict-only} & a recorded-but-uncounted ``keep'' is chain bloat
with no effect \\
relay policy & \textbf{per-identity mempool quota} & fee-less votes need a non-fee scarcity term \\
ordering at one height & \textbf{least crossing height, then least $H(r)$} & deterministic, and
first-to-quorum is the defensible priority \\
memo carriage & \textbf{coinbase field} & settled: consensus requires the ban memo
\srcintent{08-answers-round-2.md, round 5} \\
$\Delta_{\mathrm{exec}}$ & \textbf{\num{120}} (\SI{1}{\hour}) & $3\times$ the \num{40}-block consensus
reorg limit \srccode{fluxd/src/main.h}{74} \\
\SI{10}{\percent}/\SI{15}{\percent} tiers & \textbf{remove} & any destructive action there is one
party's decision (\Cref{cor:23:tiers}) \\
vote signing key & \textbf{collateral (owner) key} & weight derives from collateral; a node key is
VPS-resident (C32) \\
pool remainder & \textbf{burn pro-rata} & the only option giving voters a current-term cost
(\Cref{rem:o16-burn}) \\
$\mathsf{evidence}$ & \textbf{optional content hash} & mandatory evidence is unverifiable on chain;
absent evidence is unreviewable \\
report counts as vote & \textbf{yes, at $\hgt_0$} & a separate confirming vote from the reporter is
a redundant transaction \\
$m$ identity chamber & \textbf{\num{21}} & odd (no ties); Sybil pre-positioning costs
$(21-\kappa)\times\num{1000}\flux$ \\
$T$ identity tenure & \textbf{\num{86400}} (\SI{30}{\day}) & reuses the \PNR{} constant $K$
(\Cref{sec:pnr}); one tenure constant, two mechanisms \\
\bottomrule
\end{tabular}
\end{table}

\paragraph{The four that carry real weight.} Most of the table is mechanical. Four are not, and the
reasoning is given rather than compressed into a cell.

\textbf{$\Delta_W$ as a trailing maximum, not a mean.} The electorate-timing attack of
\Cref{sec:23:attacks} works by voting while $\vtotal^{\!\star}$ is transiently low, since the
\SI{25}{\percent} threshold is a fraction of it. A trailing maximum over seven days means the
denominator cannot be made to dip on the attacker's schedule: they must wait out the window. The
cost is that after a genuine fleet contraction the threshold stays high --- eviction becomes
temporarily harder --- and that is the correct direction for the error to point.

\textbf{Enforcement scope is the application identity, not a specification version.} If
$\mathsf{app\_hash}$ commits to one specification version, a deployer evades the ban by
re-broadcasting the same application with any trivial edit, and the quorum has to vote again for each
edit. If it extends to \emph{all} applications of the same owner, the mechanism becomes collective
punishment keyed on a field whose binding to a real party is exactly what \Cref{sec:23:measurement}
shows is weak. $(\mathsf{app\_name}, \mathsf{owner})$ is the smallest scope that is not trivially
evadable, and it matches the settled position that a banned deployer may return as a genuinely fresh
application \srcintent{08-answers-round-2.md, round 5}.

\textbf{Remove the graduated tiers.} \Cref{cor:23:tiers} is the argument: at \SI{10}{\percent} and
\SI{15}{\percent}, the measured concentration means a single party can reach the threshold alone. A
destructive action there is therefore one party's decision wearing the costume of a quorum, which is
worse than no mechanism because it carries the mechanism's legitimacy. Non-destructive actions ---
flagging, warning --- survive that objection but do no work: they neither stop harm nor inform anyone
who is not already watching the chain. One threshold, at \SI{25}{\percent}, is the honest design.

\textbf{The owner key signs, not the node key.} Voting weight is collateral-weighted, and the
collateral is controlled by the owner key; the node's \texttt{pubkey} is a VPS-resident signing key
(conflict C32). If the node key could vote, compromising $n$ virtual machines would deliver $n$
nodes' voting weight without touching a single unit of collateral --- which would make the security
of a collateral-weighted vote depend on the security of the least-hardened VPS in the fleet. The
start-transaction shape \srccode{fluxd/src/fluxnode/activefluxnode.cpp}{253--288} already carries an
owner-key signature and is the right template, which is also the shape the team described
\srcintent{08-answers-round-2.md, round 5}.

% =====================================================================
\subsection{Comparison with prior art}
\label{sec:23:priorart}
\label{sec:23:related}

\paragraph{Stake-weighted governance.} Compound and MakerDAO weight votes by token balance against
an absolute quorum expressed as a fraction of supply — the same shape as
$\vthresh\vtotal$ — and both have exhibited the plutocracy failure measured above. In Compound's
Proposal 289 (July 2024) a coordinated holder group won a majority of votes cast --- \num{682190}
COMP for against \num{633640} --- for moving \num{499000} COMP into a vault it controlled, and the
proposal was canceled before execution following community pressure rather than defeated by the
mechanism \citep{compound289}; in MakerDAO's October 2020 incident a participant used
flash-borrowed MKR, roughly \$7M worth, to pass a whitelisting vote and returned the tokens in the
same transaction, after which Maker extended its governance security module delay from
\SI{12}{\hour} to \SI{72}{\hour} \citep{makerflashloan2020}. \emph{What transfers:} that stake-weighting
concentrates decision power in proportion to capital; that delegation concentrates it further; and
that a post-decision delay — Compound's timelock, Maker's pause — is the cheapest available safety
margin, which here is $\Delta_{\mathrm{exec}}$. \emph{What does not transfer:} flash-loan capture.
Voting weight here would be \emph{locked} collateral with 100-block maturity and a confirmation
requirement, and the electorate would be frozen at $\hgt_0$, so weight cannot be borrowed into a
vote after the report exists. That is a genuine advantage of collateral weighting and the paper
claims it. The plutocracy failure, by contrast, transfers exactly. In both incidents the decisive
set was small — a coordinated holder group in one case, a single flash-borrowing participant in the
other — which is the same shape as $\kappa(0.25) = 4$, with one difference that cuts against Flux:
here the decisive set is a standing property of the collateral distribution rather than something
assembled for one vote.

\paragraph{Sybil resistance.} \citet{douceur2002sybil} establishes that without a trusted identity
authority or a resource cost, a single party can present arbitrarily many identities — which is
exactly why the settled weighting binds voting power to locked collateral and why every
concentration remedy that decouples them (weight caps, per-identity quadratic weighting) is defeated
at zero cost (\Cref{sec:23:negative}). Personhood-based approaches do not transfer: Flux operators
are businesses and machines, and no personhood attestation is available to a \FluxOS{} node. What
does transfer is the \emph{bicameral} shape used where token weight is paired with an
attested-identity house; \Cref{def:23:twochamber} is that structure with tenure standing in for
personhood — a much weaker identity, which is precisely why it claims detectability rather than
resistance.

\paragraph{Content moderation in peer-to-peer systems.} \citet{benet2014ipfs} describes a
content-addressed store; the surrounding ecosystem publishes an opt-in denylist of hashed
identifiers that each node \emph{may} apply. Freenet has no removal at all by design; the fediverse
moderates per instance by defederation; Tor exit policies and relay policies are per operator.
\emph{What transfers:} the per-node denylist is exactly the shape of Flux's deployed blocklist —
except that Flux's is \emph{mandatory} on every node and centrally authored, which makes Flux more
centralised than IPFS in this one respect. \emph{What does not:} every system above stores or relays
\emph{bytes}. Flux runs \emph{processes} with network egress that can phish, command, spam and
attack third parties, which is why a network-wide removal — stronger than anything those systems do
— is on the table at all. Note that the intended scope is broader still: by
\Cref{rem:content-veto} the quorum is a content veto, so Flux is choosing a power that none of the
comparison systems exercises, and choosing it knowingly. It is also why concentration matters more here: a denylist a node may ignore is a
recommendation; a ban carried in the chain that every node must honour is a power.

\paragraph{Filecoin, Arweave, Akash.} \citet{filecoin2017} specifies a market in which storage
providers choose which deals to accept; the protocol itself removes nothing. Arweave nodes run local
content policies over what they store and serve, and the protocol cannot remove. Akash providers set
per-provider acceptance attributes, and no network-wide eviction exists \citep{akash}. \emph{What
transfers:} the position that the protocol should not hold a content veto — which is the roadmap's
own principle and which the quorum contradicts unless the taxonomy is confined to behaviour.
\emph{What does not:} their tenants' data is inert and Flux's tenants' processes are not, and a
provider-level ``I will not host this'' is categorically different from a network-level ``nobody may
host this''. The quorum is the latter. The paper presents it as a deliberate departure from the
Filecoin and Arweave position, justified by the difference between storing and executing, rather
than as an industry norm.

% =====================================================================
\subsection{Open parameters}
\label{sec:23:open}

The on-chain shape settles six things the earlier specification carried as parameters, and it is
worth listing them as gains: the voting window ($\Delta_{\mathrm{vote}} = \num{2880}$ blocks), the
per-application cooldown after an unmet report (zero), reporter eligibility
($\reporters = I(\hgt_0)$, by \Cref{def:23:capacity}), the network-wide cap on concurrently open
reports ($\lvert N(\hgt) \rvert$, by \Cref{cor:23:capacity}), the entire bond question, which
disappears with the bond, and producer obligation --- ban-memo inclusion, over the first unbanned
carried report in the settled order, is consensus-required rather than discretionary
(\Cref{prop:23:inclusion-required}). Eleven
further elements were carried as open parameters through the earlier drafts, each marked in place
above with its domain and its trade-off, and all but the first are now settled at the values of
\Cref{tab:23-recommended}: the category taxonomy, which remains open as an enumeration task
(\Cref{rem:content-veto}); $\Delta_W$; what $\mathsf{app\_hash}$ hashes, which is
also the enforcement scope; the $\mathsf{decision}$ field and the relay policy for fee-less
transactions; the ordering rule when two reports cross in
one slot, and how the memo is carried in the coinbase; $\Delta_{\mathrm{exec}}$; the actions at the
graduated \SI{10}{\percent} and \SI{15}{\percent} tiers; which transaction shape the vote follows
and hence which key signs it; the report's evidence field and whether a report is itself a vote;
$(m, T)$ for the identity chamber; and the disposition of the pool's unearned remainder.

Three further items are recorded here rather than marked in place, because they are not parameters of
this mechanism. Whether opt-in delegation is adopted, and with what revocation semantics
(\Cref{sec:23:remedies}). Whether \FDM{} and \PDNS{} honour bans by dropping routes at
$\hgt_b + \Delta_{\mathrm{exec}}$ — without which a banned application would remain reachable at its
name until the last node had removed it, and with which the Foundation-operated reachability plane
acquires a new dependency on moderation state. And the sequencing question this section
inherits from \Cref{sec:23:endgame}: signing the deployed policy documents is strictly cheaper than
the quorum, requires no consensus change, and should not wait for it. Whether the ban set is ever
pruned is no longer open: \Cref{prop:23:monotone} and \Cref{op:23:reinstatement} together settle it
as no --- a ban is final, and the ban set only grows.

Three items cannot be set by parameter choice. Two are \emph{verifications} in \fluxd{}: whether
the emergency-block path is subject to the coinbase rules the ban memo would rely on --- now checked
against source and answered favourably (\Cref{rem:emergency-coinbase}) --- and whether re-addressing
resets \texttt{added\_height}, which decides whether the tenure gate of \Cref{def:23:twochamber} does
any work at all and remains unverified. The third cannot be resolved at all from public data:
beneficial ownership behind the payment addresses that carry the voting weight
(\Cref{conj:23:custody}). The first two are prerequisites, not refinements. The third is the reason
$\kappa(0.25) = 4$ is printed as a measurement of \emph{identities} and never as a count of
\emph{people}.

%% Drain this section's float queue before the next begins.
\FloatBarrier
